\chapter{用观察性研究估计因果效应}
\chapterauthors{Thomas B. Newman、Warren S. Browner}
\begingroup\setlength{\tabcolsep}{3pt}
\sourcepage{276}
多数观察性研究旨在估计某种暴露对结局的因果效应，例如食用红肉对结肠癌风险的影响。（以预测为目的的观察性研究属于例外，如第13章讨论的诊断与预后检查研究。）如第1章所述，这一过程首先是从研究样本中测得的变量，推断目标人群中所关注的现象。但到目前为止，我们讨论的推断都只涉及目标人群中的关联。要估计\term{因果效应}（causal effects），还需进一步工作。因果效应比关联更值得关注，因为它有助于理解疾病背后的病理生理机制，找出预防疾病发生的方法，甚至提示可能的治疗方法。

\section{理解因果关系的反事实框架}
为了理解因果效应，可以做一个思想实验：比较目标人群中\emph{所有人均暴露}（整个目标人群都吃红肉）时的结局（结肠癌）风险，与\emph{同一人群无人暴露}（没有人吃红肉）时的风险。\footnote{为简化讨论，我们把预测因素与结局都作为二分类变量；一般因果效应的估计同样适用这些原则。}当然，在现实世界中，每个人不是已经暴露，就是没有暴露。因此，我们无法知道这些所谓的\term{反事实}（counterfactual）世界会发生什么；在这些世界中，许多人的暴露状态都“与事实相反”。最适合估计因果效应的研究，应使研究者能够估计并比较：在这些反事实暴露状态下，本来会出现怎样的反事实结局（图10.1）。
\providecommand{\chTenCounterfactual}{\begin{tikzpicture}[font=\small,>=Stealth]
\fill[dcrlight] (-5,0) rectangle (5,4.2);\fill[dcrgold!15] (-5,-4.2) rectangle (5,0);
\draw[gray] (0,-4.2)--(0,4.2);\draw[gray] (-5,0)--(5,0);
\node[text=dcrblue] at (0,4.6) {现实世界};\node[text=dcrgold] at (0,-4.6) {反事实世界};
\node at (-2.5,3.8) {有暴露};\node at (2.5,3.8) {无暴露};
\node[anchor=west] at (-4.9,.3) {A};\node[anchor=east] at (4.9,.3) {B};
\node[anchor=west] at (-4.9,-.3) {C};\node[anchor=east] at (4.9,-.3) {D};
\node[draw,fill=dcrblue!20] at (-3.5,2.7) {男＋};\node[draw,fill=yellow!40] at (-1.4,1.5) {女＋};
\node[draw,fill=gray!30] at (1.2,2.8) {男};\node[draw,fill=dcrpurple!30] at (2.6,2) {女};\node[draw,fill=dcrteal!25] at (4,1.2) {女};
\node[draw,fill=gray!30] at (-4,-1.2) {男＋};\node[draw,fill=dcrpurple!30] at (-2.6,-2) {女＋};\node[draw,fill=dcrteal!25] at (-1.2,-2.8) {女＋};
\node[draw,fill=dcrblue!20] at (1.4,-1.5) {男};\node[draw,fill=yellow!40] at (3.5,-2.7) {女};
\draw[->,line width=2pt] (-.9,.9)--(.9,-.9);\draw[->,line width=2pt,gray] (.9,.9)--(-.9,-.9);
\node[circle,draw,fill=white,text width=22mm,align=center,minimum size=28mm] at (0,0) {反事实\\转换器};
\end{tikzpicture}}
\providecommand{\chTenPlan}{\begin{tikzpicture}[font=\small,>=Stealth]
\node[draw=dcrblue,text width=48mm,align=left,minimum height=65mm] (q) at (0,0) {\textbf{研究问题}\\[6pt]\textbf{目标人群}\\全部成人\\[9pt]\textbf{目标现象}\\原因：红肉摄入\\$\downarrow$ 因果效应\\效应：结肠癌};
\node[draw=dcrblue,text width=48mm,align=left,minimum height=65mm] (p) at (8,0) {\textbf{研究计划}\\[6pt]\textbf{预定样本}\\符合选择标准的住院\\结肠癌病例和对照\\[9pt]\textbf{预定变量}\\预测：食物频率问卷的红肉摄入\\$\updownarrow$ 关联\\结局：病历结肠癌诊断};
\draw[->,thick,dcrblue] (q)--node[above]{设计}(p);
\node[draw=dcrred,minimum width=48mm,minimum height=12mm] (u) at (0,-4.7) {当前认识};
\node[draw=dcrred,minimum width=48mm,minimum height=12mm] (c) at (8,-4.7) {研究结论};
\draw[->,thick,dcrred] (c)--node[above]{推断}(u);
\end{tikzpicture}}
\providecommand{\chTenColliders}{\begin{tikzpicture}[font=\small,>=Stealth]
\foreach \yy/\aa/\bb/\cc/\panel in {0/节食/癌症/体重下降/A,-3/父母吸烟/早产/低出生体重/B}{\begin{scope}[yshift=\yy cm]
\node at (-3,.3) {\panel};\node[text=dcrpurple] (x) at (-1.5,.5) {\aa};\node[text=dcrgold] (y) at (2.2,.5) {\bb};
\node[draw=dcrred,text=dcrteal] (z) at (.4,-.8) {\cc};\draw[->,dcrpurple] (x)--(z);\draw[->,dcrgold] (y)--(z);
\end{scope}}\end{tikzpicture}}
\providecommand{\chTenTableModel}[1]{\begin{tikzpicture}[font=\footnotesize,>=Stealth,thick,baseline=(current bounding box.north)]
\ifcase#1
\or
\node (m) at (0,1.2) {吃红肉};\node (v) at (0,.3) {CVD};\node[draw] (h) at (2.2,-.5) {住院};\node (c) at (0,-1.3) {CRC};
\draw[->,dcrblue] (m)--(v);\draw[->,dcrblue] (v)--(h);\draw[->,dcrblue] (c)--(h);
\or
\node (c) at (0,1) {CRC};\node (a) at (0,0) {贫血};\node (m) at (0,-1) {吃红肉};\draw[->,dcrblue] (c)--(a);\draw[->,dcrblue] (a)--(m);
\or
\node (s) at (0,1.2) {男性性别};\node (m) at (-1.1,0) {吃红肉};\node (c) at (1.5,0) {CRC};\draw[->,dcrblue] (s)--(m);\draw[->,dcrblue] (s)--(c);\draw[->,dcrblue] (m)--(c);
\fi\end{tikzpicture}}
\providecommand{\chTenSexDAG}{\begin{tikzpicture}[font=\small,>=Stealth,thick]
\node (s) at (0,1.4) {性别};\node (m) at (-1.2,0) {吃红肉};\node (c) at (1.5,0) {CRC};
\draw[->,dcrblue] (s)--(m);\draw[->,dcrblue] (s)--(c);\draw[->,dcrblue] (m)--(c);\end{tikzpicture}}
\providecommand{\chTenTablePath}[1]{\begin{tikzpicture}[font=\scriptsize,>=Stealth,thick,x=.55cm,y=.55cm,baseline=(current bounding box.north)]
\ifcase#1
\or
\node (p) at (0,0) {LDL};\node (o) at (4,0) {卒中};\node[align=center] (z) at (2,-1.7) {颈动脉\\狭窄};\draw[->] (p)--(o);\draw[->,dcrblue] (p)--(z);\draw[->,dcrblue] (z)--(o);
\or
\node (p) at (0,0) {吸大麻};\node (o) at (4,0) {肺癌};\node (z) at (2,1.7) {吸烟};\draw[->] (p)--(o);\draw[->,dcrblue] (z)--(p);\draw[->,dcrblue] (z)--(o);
\or
\node[align=center] (p) at (0,0) {低热量\\饮食};\node (o) at (4,0) {胃癌};\node[align=center] (z) at (2,-1.7) {体重下降};\draw[->] (p)--(o);\draw[->,dcrblue] (p)--(z);\draw[->,dcrblue] (o)--(z);
\fi\end{tikzpicture}}
\providecommand{\chTenDAG}[1]{\begin{tikzpicture}[font=\small,>=Stealth,thick]
\ifcase#1
\or\node[text=dcrpurple](p)at(0,0){头部创伤};\node[text=dcrgold](o)at(7,0){帕金森病};\node[text=dcrteal](z)at(3.5,1.7){骑车戴头盔};\draw[->,dcrblue](z)--(p);\draw[->](p)--(o);
\or\node[text=dcrpurple](p)at(0,0){A型血};\node[text=dcrgold](o)at(7,0){因COVID-19住院};\node[text=dcrteal](z)at(3.5,1.7){男性};\draw[->,dcrblue](z)--(o);\draw[->](p)--(o);
\or\node[text=dcrpurple](p)at(0,0){雄激素水平};\node[text=dcrgold](o)at(7,0){肾癌};\node[text=dcrteal](z)at(3.5,-1.7){男性型脱发};\draw[->,dcrblue](p)--(z);\draw[->](p)--(o);
\or\node[text=dcrpurple](p)at(0,0){终生铅暴露};\node[text=dcrgold](o)at(7,0){退行性关节病};\node[text=dcrteal](z)at(3.5,-1.7){使用镇痛药};\draw[->,dcrblue](o)--(z);\draw[->](p)--(o);
\or\node[text=dcrpurple](p)at(0,0){LDL胆固醇};\node[text=dcrgold](o)at(7,0){卒中};\node[text=dcrteal](z)at(3.5,-1.7){颈动脉狭窄};\draw[->,dcrblue](p)--(z);\draw[->,dcrblue](z)--(o);\draw[->](p)--(o);
\or\node[text=dcrpurple](p)at(0,0){低热量饮食};\node[text=dcrgold](o)at(7,0){胃癌};\node[text=dcrteal](z)at(3.5,-1.7){体重下降};\draw[->,dcrblue](p)--(z);\draw[->,dcrblue](o)--(z);\draw[->](p)--(o);
\or\node[text=dcrpurple](p)at(0,0){吸大麻};\node[text=dcrgold](o)at(7,0){肺癌};\node[text=dcrteal](z)at(3.5,1.7){吸烟};\draw[->,dcrblue](z)--(p);\draw[->,dcrblue](z)--(o);\draw[->](p)--(o);
\or\node[text=dcrpurple](p)at(0,0){肥胖};\node[text=dcrgold](o)at(7,0){关节炎};\node[text=dcrteal](z)at(3.5,-1.7){久坐生活方式};\draw[->,dcrblue](o)--(z);\draw[->,dcrblue](z)--(p);\draw[->](p)--(o);
\or\node[text=dcrpurple](p)at(0,0){激素替代治疗};\node[text=dcrgold](o)at(7,0){子宫内膜癌};\node[text=dcrteal](z)at(3.5,-1.7){绝经后出血};\draw[->,dcrblue](p)--(z);\draw[->,dcrblue](o)--(z);\draw[->](p)--(o);
\or\node[text=dcrpurple](p)at(0,0){吃红肉};\node[text=dcrgold](o)at(7,0){结直肠癌};\node[text=dcrteal](z)at(3,2.3){健康生活方式};\node[text=dcrteal](s)at(5.4,1.2){癌症筛查};\draw[->,dcrblue](z)--(p);\draw[->,dcrblue](z)--(s);\draw[->,dcrblue](s)--(o);\draw[->](p)--(o);
\fi\end{tikzpicture}}
\providecommand{\chTenManyDAG}{\begin{tikzpicture}[font=\small,>=Stealth,thick]
\node[text=dcrpurple](p)at(0,0){\underline{红肉摄入}};\node[text=dcrgold](o)at(8,0){\underline{结直肠癌}};
\node[text=dcrteal](pol)at(-.5,2){政治观点};\node[text=dcrteal](f)at(3.8,2){\underline{膳食纤维}};\node[text=dcrteal](g)at(8,2){遗传多态性};
\node[text=dcrteal](m)at(4,-1.6){微生物组};\draw[dotted,dcrteal] (3.3,-1.85)--(4.7,-1.85);
\node[text=dcrteal](b)at(4,-3.2){粪便隐血试验阳性};
\draw[->](p)--(o);\foreach \a/\b in {pol/p,f/p,f/o,g/o,p/m,m/o,p/b,o/b}{\draw[->,dcrblue](\a)--(\b);}
\end{tikzpicture}}
\providecommand{\chTenIQ}{\begin{tikzpicture}[font=\small,>=Stealth,x=.85cm,y=.18cm]
\draw[->] (6,90)--(16.5,90);\node at(11,85.5){父母受教育年限};\draw[->] (6,90)--(6,122) node[above] {儿童IQ};
\foreach \x in {6,8,10,12,14,16}{\draw(\x,90)--(\x,89.5) node[below]{\x};}
\foreach \y in {90,95,100,105,110,115,120}{\draw(6,\y)--(5.9,\y) node[left]{\y};}
\draw[dashed,dcrgold,thick](6,95)--(15.8,119.5);
\foreach \y in {95,105,110}{\draw[gray,thin](6,\y)--(16,\y);}
\fill[dcrblue](12,110)circle(2pt) node[below right]{$A$};\fill[dcrblue](10,105)circle(2pt) node[below right]{$A'$};\fill[dcrblue](10,95)circle(2pt) node[below right]{$B$};
\node[align=center,anchor=east] (a)at(11.5,117){平均IQ\\正常铅水平};\draw[->](a)--(11.9,110.4);
\node[align=center,anchor=east](ap)at(9.8,112){正常铅平均IQ\\按父母教育调整};\draw[->](ap)--(10,105.4);
\node[align=center](b)at(8.7,92.5){平均IQ\\高铅水平};\draw[->](b)--(9.85,94.8);
\draw[<->](14.1,105)--(14.1,110);\node[align=left,anchor=west,text width=21mm]at(14.2,107.5){父母教育\\所致IQ差};
\draw[<->](14.1,95)--(14.1,105);\node[align=left,anchor=west,text width=21mm]at(14.2,100){铅所致\\IQ差};
\end{tikzpicture}}

\sourcepage{277}
\par
\begin{center}\begin{minipage}{\textwidth}\centering\chTenCounterfactual
\captionof*{figure}{\textbf{图10.1\quad 用反事实世界理解因果效应的思想实验。}在现实世界中，通过比较吃红肉者（A）与另外一些不吃红肉者（B）的结局（如结肠癌），估计暴露（吃红肉）的效应，并试图控制两组的其他差异，如图中女性较多不吃红肉。在思想实验中，反事实转换器创造出与现实人物完全相同、但暴露相反的人。所有人都吃红肉的世界（A＋C），与所有人都不吃红肉、其他方面完全相同的世界（B＋D）之间，结肠癌率的差，即红肉对结肠癌的因果效应。图中“＋”表示吃红肉。}
\end{minipage}\end{center}
\par

最简单的方法是随机试验，下章将作详细讨论。这里先关注试验如何用于估计反事实结局。首先，需要招募因果效应对其有意义的人，也就是至少在理论上可以处于任一种暴露水平的人。例如，有些人可能由于宗教、伦理或医学原因，在任何情况下都不能或不愿吃红肉，我们就应将其排除在试验之外。对于坚持一定要吃红肉的人，也应如此。\sourcepage{278}然后，把其余参与者随机分配到不同暴露状态，形成相似的两组，用流行病学术语说，就是“\term{可交换}”（exchangeable）的两组。其中一组暴露（本例中分配为吃红肉），另一组不暴露（分配为不吃红肉）。随机化使两组可交换，因此吃红肉组的结肠癌率，可用于估计不吃红肉组假如吃了红肉会发生什么；同样，不吃红肉组的结局，可用于估计吃红肉组假如不吃红肉会怎样。因此，随机试验中观察到的组间风险差，在参与者良好依从治疗分配的条件下，就是因果效应的有效估计。在观察性研究中增强因果推断，通常要尽可能模拟一项回答该研究问题的随机试验。

\section{观察性研究关联为何可能不同于因果效应}
不过，在暴露不受研究者控制的观察性研究中，试图模拟随机试验可能很有挑战。再考虑吃红肉是否引起结肠癌的问题。对这一主题开展观察性研究很有意义：结肠癌是美国癌症死亡的第二大原因（1）；对一些人来说，对哺乳动物的同情，以及红肉生产已经证实的环境不利影响，尤其是牛羊等反刍动物肉类的生产（2），还不足以促使他们改变饮食；而将人随机分配为吃或不吃红肉，再随访足够长时间以观察结肠癌结局，又极不现实。让我们将红肉暴露作为一个二分类变量（是／否），并假定一项观察性研究发现，吃红肉者发生结肠癌的风险约为不吃红肉者的两倍。

一种可能性，想必也是研究者最关心的可能性，是吃红肉使结肠癌风险加倍。但在下结论之前，必须考虑四种可能使这一结论不成立的原因（表10.1）：偶然、偏倚、混杂和果因。\footnote{有些流行病学家把混杂视为一种偏倚，因为它和所有偏倚一样，会使研究所估计的参数（即因果效应的度量）失真。我们倾向于单独讨论混杂，因为其他类型偏倚会使研究中的关联不同于目标人群中的真实关联，无论目标是否为估计因果效应，都有问题。相比之下，混杂只有在目标是估计因果效应时才成问题：它使目标人群中的关联不同于因果效应。}要准确估计吃红肉对结肠癌的因果效应，就须减少这四种可能性造成的失真。
\sourcepage{279}
\begin{longtable}{L{23mm}L{28mm}L{42mm}L{56mm}}
\caption*{\textbf{表10.1\quad 红肉摄入与结直肠癌研究中因果效应估计失真的四个原因}}\\\toprule
原因&研究估计反映目标人群关联吗？&人群中实际情况&因果模型\\\midrule\endfirsthead
\toprule 原因&反映人群关联吗？&实际情况&因果模型\\\midrule\endhead
1. 偶然&否&随机误差使关联失真&吃红肉与CRC之间没有连接箭头\\
2. 偏倚&否&系统误差使关联失真&医院病例对照研究中的选择偏倚（见正文）\par\medskip\chTenTableModel{1}\\
3. 果因&是&结肠癌使人吃更多红肉&\chTenTableModel{2}\\
4. 混杂&是&第三因素同时引起吃红肉与结肠癌&\chTenTableModel{3}\\\bottomrule
\multicolumn{4}{p{154mm}}{\footnotesize CVD：心血管疾病；CRC：结直肠癌。}\\
\end{longtable}
\sourcepage{280}
\par
\begin{figure}[htbp]\centering\chTenPlan
\caption*{\textbf{图10.2\quad 考虑研究问题与研究计划的差异，以尽量减少偏倚。}}
\end{figure}
\par

如果解释是偶然或偏倚，那么在研究中，吃红肉者的结肠癌风险会加倍，尽管人群中并不存在这一关联。表面上的风险加倍，实际上源于运气不好，或研究设计、实施、解释中存在问题。

另外两种解释，即果因和混杂，都是真实的生物学现象，因此样本中的关联可能反映人群中的关联。不过，这种关联可能并不完全是暴露的因果效应。在果因情形中，结肠癌使人吃更多红肉，例如肠道失血造成贫血后，患者可能渴望吃红肉。这只是因果方向倒转，主要是回顾性研究需要考虑的问题。最后一种可能性是混杂：男性性别认同（下文简称“男性性别”）等第三因素，同时引起红肉摄入和结肠癌。

本章接下来介绍，在观察性研究中估计因果效应时，如何估计并尽量减少上述四种错误造成的失真。这些策略既可用于研究设计阶段，也可用于结果分析阶段。\sourcepage{281}本书虽侧重研究设计，但了解可选的分析方法会影响设计选择，因此两个方面都将讨论。

\section{尽量减少偶然所致错误}
假设目标人群中50\%的人吃红肉，但红肉摄入与结肠癌没有关联。如果随机抽取20名结肠癌病例和20名对照，预期每组约有10人吃红肉，即20人的50\%。但是，仅仅由于偶然，20名病例中可能纳入了14名红肉食用者，而20名对照中只有6名。如果发生这种情况，研究就会观察到红肉摄入与结肠癌之间虚假、却有统计学意义的关联：比值比为5.4，95\%置信区间1.2～26，$P=0.03$。

偶然有时称为随机误差，因为它没有背后的解释。由随机误差造成的关联，如果有统计学意义，就称为I类错误（第5章）。研究设计和分析阶段都有减少随机误差的策略（表10.2）。提高测量精密度和增加样本量等设计策略，分别见第4章和第6章。分析阶段计算$P$值，可帮助研究者将观察到的关联与仅凭偶然可能出现的关联相比较，量化关联的大小。例如，上段的$P=0.03$意味着，目标人群中不存在关联时，仅凭偶然，约3\%的情况下也会出现至少这么强的关联。

\begin{longtable}{L{25mm}L{65mm}L{59mm}}
\caption*{\textbf{表10.2\quad 正确估计目标人群关联，减少因果效应估计失真}}\\\toprule
原因&设计阶段：预防失真&分析阶段：评估失真\\\midrule\endfirsthead
\toprule 原因&设计阶段&分析阶段\\\midrule\endhead
偶然（随机误差）&增大样本及其他提高精度的方法（第4、6章）&计算$P$值和置信区间，结合既有结果解释（第5、11章）\\
偏倚（系统误差）&仔细考虑问题与计划的每项差异后果（图10.2），必要时修改计划&检查与其他、尤其不同设计研究是否一致\\
&不用受暴露影响的变量作纳入或匹配条件&不控制受预测变量影响的变量\\
&设盲及其他方法，尽量减少暴露组间测量准确度、精密度差异&寻找观察者、中心之间均数、标准差的系统差异\\
&收集可评估偏倚程度的额外资料（证伪检验），如预计不与暴露相关的结局，反之亦然&分析额外资料（证伪检验），看是否提示潜在偏倚已发生\\\bottomrule
\end{longtable}
\sourcepage{282}
置信区间比$P$值更有用，它显示描述关联的参数（风险比、比值比等），在本研究估计的随机误差范围内可能取哪些值。对于没有统计学意义的结果，置信区间尤其有帮助，因为它能显示是否可能由于偶然，而漏掉了一个大且重要的效应（3、4）。

\section{尽量减少偏倚所致错误}
偏倚也称系统误差，种类很多，如何处理它们是本书的重要主题。除了第3、4、8、9章介绍的具体策略，这里再补充一种降低偏倚可能性的一般思路。

如第1章所述，原始研究问题与研究实际回答的问题，几乎总有差异。这些差异反映了为使研究可行而作出的妥协，以及研究设计或实施中的错误。如果这些差异使研究所得估计系统性地偏离原本要估计的量，就会出现偏倚。如第9章所述，偏倚主要分为\term{选择偏倚}（selection bias，亦称抽样偏倚）和\term{测量偏倚}（measurement bias）两类。在设计和分析阶段，都有尽量减少偏倚的策略（表10.2）。

\sourcepage{283}
\subsection{设计阶段}
像图10.2那样把研究问题放在研究计划旁，逐一思考：
\begin{enumerate}
\item 参与者样本（病例／对照或暴露／未暴露）代表关注人群吗？
\item 暴露测量代表关注的预测因素吗？
\item 结局测量代表关注的结局吗？
\end{enumerate}
对于每一个回答为“不”或“可能不”的问题，都要考虑偏倚影响的是其中一组还是两组（例如病例与对照，或暴露与未暴露参与者），对两组的影响是否相似，以及偏倚是否可能大到足以改变研究问题的答案。

用红肉与结肠癌的例子说明这一点。考虑一项医院病例对照研究，对照从因结肠癌以外疾病住院的患者中抽取。如果其中许多人患有由肉类摄入引起的心血管疾病，对照样本就不能代表结肠癌病例所来自的目标人群：对照组中吃肉者过多，会使研究中的红肉与结肠癌关联低于人群中的关联，甚至可能发现吃红肉具有预防结肠癌的作用。这就是选择偏倚。\footnote{附录10A讨论碰撞点分层偏倚和DAG时，将再次讨论这一例子。}

测量偏倚也是这项研究的潜在问题。例如，如果研究使用既有资料，而红肉摄入量是用24小时食物频率问卷测量的，会怎样？参与者回忆的24小时摄入量，可能不同于当天实际红肉摄入量，也可能不同于过去或未来通常的摄入量。差异可能仅由偶然造成，例如恰巧抽到他吃鸡块的那一天；也可能由系统误差造成，例如人们一般倾向于少报自己认为不健康的行为。

测量偏倚还可能影响研究结局的确定方式。多数队列研究把结局发生估计为一个比例或率，分子是随访期间新出现的结局数。为正确估计，必须确保各暴露组以同样方式确定结局，例如使负责判定参与者是否出现结局的人，不知道其暴露状态。\sourcepage{284}分母则是处于风险中的人群（累积发病率）或风险人时（发病密度）。处理这些分母可能很棘手，尤其是随访期间不止一次测量暴露时。如第8章所述，将观察性研究设计为模拟“目标试验”，有助于避免不死时间偏倚，以及把实际上未处于风险中的人计作处于风险，或反过来的其他错误（5—7）。

下一步是考虑每种潜在偏倚的预防策略，例如在病例对照研究中选择不止一个对照组（第9章），或采用第4章介绍的减少测量偏倚的方法。对于每种情况，都须判断偏倚发生的可能性，以及通过修改研究计划预防它有多容易。如果容易预防，就修改计划，再问一遍上述三个问题。如果不容易预防，就根据偏倚发生的可能性，以及它会使所估计关联失真到什么程度，决定研究是否仍值得开展。

潜在偏倚可能无法避免，也可能预防成本很高，或者研究者尚不清楚它会造成多大问题。遇到这些情况，应考虑在研究设计中加入额外资料的收集，以评估偏倚的严重程度。这类检验称为\term{证伪检验}（falsification tests），应当预先规定并写入研究设计（8）。例如，第9章介绍的恶性黑色素瘤与日晒的双生子病例对照研究中，研究者担心回忆偏倚，便询问一对双生子中的双方：童年时谁做日光浴更多。患黑色素瘤者报告自己的日晒更多，而作为对照的双生子则报告两人的日晒相近。虽然无法知道谁的回答正确，答案的差异提示童年日晒的回忆存在问题。\footnote{无法确定这是病例相较对照高估暴露的回忆偏倚，还是不完整回忆造成的测量误差，而这种误差对病例和对照的影响相似。}如果担心24小时食物频率问卷不能准确反映通常的红肉摄入量，可以在部分病例和对照中加用更详细、但更耗时的7天饮食日记，以了解它与24小时食物频率问卷的一致程度。同样，如果担心红肉并不引起结肠癌，而是提高结肠癌患者的生存率，从而使红肉食用者在结肠癌存活病例中所占比例过高，病例对照研究就可以找到已经死亡的结肠癌患者，\sourcepage{285}再访谈其尚存的伴侣，询问患者过去的饮食。

\subsection{分析阶段}
资料收集完成后，目标就从尽量减少偏倚转为评估偏倚可能有多严重。第一步是分析为预先规定的证伪检验所收集的资料。例如，研究者预料参与者可能无法准确回忆红肉摄入量，便可预先询问病例和对照对自己的答案有多大把握。分析时再按其对红肉摄入量的确信程度分层，考察对暴露史更有把握的人中，红肉与结肠癌的关联是否更强。

研究者还可以查看其他研究的结果。如果结论一致，关联由偏倚造成的可能性就较低；如果其他研究采用了不同设计，因此不太可能具有相同的偏倚，这一点就更有说服力。应当多积极地寻求额外信息，以及报告研究时如何讨论这些问题，都需要判断；征求同事的意见会有帮助。

\section{辨清方向：果因}
偶然和偏倚可以解释，为什么样本中测量变量的关联估计不一定反映目标人群中所关注现象的真实关联。但是，即使关联估计有效，它也不一定代表因果效应。

一种可能是本末倒置：结局引起了预测因素（表10.3）。果因常是横断面研究和病例对照研究中的问题：是久坐生活方式引起肥胖，还是反过来？病例交叉研究也可能遇到果因。例如，第9章介绍的手机使用与机动车事故研究（9）中，准确知道通话和车祸发生的时间很重要，因为可能是车祸使驾驶者打手机报告事故，而不是驾驶者分心导致了车祸。

\begin{table}[htbp]\centering\small
\caption*{\textbf{表10.3\quad 观察性研究中辨清因果方向的策略}}
\begin{tabularx}{\textwidth}{Y Y}\toprule
设计阶段&分析阶段\\\midrule
纵向研究以知谁先发生；取得变量历史顺序资料；最终解决方法是随机试验&考虑生物学合理性；比较预测因素暴露后立即与较晚时的关联强度；参考其他设计研究结果\\\bottomrule
\end{tabularx}
\end{table}

\sourcepage{286}
研究疾病病因的队列研究较少遇到果因问题，因为可以在尚未患病的参与者中测量危险因素。但即便是队列研究，如果疾病潜伏期很长，而且基线时无法识别已有亚临床疾病者，仍可能发生果因。例如，2型糖尿病与随后患胰腺癌的风险有关。其中一部分关联很可能是果因，因为胰腺癌可能损伤分泌胰岛素的胰岛细胞，从而引起糖尿病。胰腺癌风险在糖尿病确诊后不久最高，这符合果因的解释（10）。关联随糖尿病病程延长而减弱，但在糖尿病发生4年甚至更久后，胰腺癌仍有超额风险（10—12），提示其中至少一部分关系可能是因果。

这一例子说明了排除果因的一般方法：考察假定的原因与效应之间的时间延长时，关联是否减弱。当然，对于很快产生作用的短期暴露，因果效应本身也会随时间减弱。因此，这种策略最适合糖尿病等慢性暴露，因为其累积暴露随时间增加；如果存在因果关系，预期关联应随时间增强，而不是减弱。

第二种方法，是比较果因与因果的生物学合理性。在这一例子中，胰腺癌可能损伤胰腺，所以果因在生物学上合理。不过，糖尿病超过10年不仅与胰腺癌风险升高有关，也与多种其他癌症风险升高有关（12），这一观察增强了糖尿病引起胰腺癌的生物学合理性。

\section{混杂}
人群中的关联可能不同于因果效应，另一个原因是\term{混杂}。如前所述，男性性别等第三因素同时引起红肉摄入和结肠癌时，就会发生混杂。\footnote{这里略作简化：混杂因素也可以不引起暴露，而是与暴露具有共同原因。例如，筛查不使人少吃红肉，但二者可能同受健康意识影响：健康意识高者更依从筛查，也少吃红肉。这时应控制筛查，防止它混杂红肉对结肠癌发生的作用。要更充分地理解混杂，需要使用DAG，见附录10A。}\sourcepage{287}担心男性性别造成混杂是合理的，因为吃红肉与一些男性对男子气概的认识有关（13），而且男性在各年龄的结肠癌发病率更高（14）。如果确实如此，目标人群中红肉摄入与结肠癌的关联，就可能不完全代表因果效应。附录10B给出数字例子，说明性别混杂如何部分造成红肉与结肠癌的关联。

如果吃红肉改变了肠道微生物组，随后由微生物组引起结肠癌，又会怎样（15）？这时，微生物组是吃红肉对结肠癌因果效应的\term{中介因素}（mediator），而不是混杂因素。一般来说，本章后面介绍的混杂应对策略，不应用于预测因素到结局因果路径上的因素，也就是中介因素。

除偏倚以外，混杂往往是因果关系唯一可能的替代解释，也是最需要设法排除的解释。它同时又最难排除，因此本章余下的大部分篇幅，都用于讨论处理混杂因素的策略。不过，这些策略都需要判断；再精巧的流行病学或统计学方法，也不能替代对背后生物学的理解。

\section{设计阶段应对混杂}
多数应对混杂变量的策略，都要求研究者测量这些变量。因此，可以先列出可能引起暴露（或与暴露具有共同原因）、同时又引起结局的变量，例如年龄和性别。随后，研究者必须从设计和分析策略中作出选择，控制这些潜在混杂变量的影响。

前两项设计阶段策略（表10.4），即\term{限定}（specification）和匹配，都涉及改变抽样方式。病例对照研究可使病例和对照在混杂变量上具有可比的值；队列研究则可使暴露者与未暴露者具有可比的值。这样，就排除了该混杂因素对预测因素与结局之间任何观察关联的解释。第三项策略是利用机会性研究设计，它只适用于具备合适条件的某些研究问题。不过，在适用时，这类设计不仅能减少或消除已测量变量造成的混杂，也能减少或消除未测量变量造成的混杂，其能力可与随机试验相似。
\sourcepage{288}
\begin{longtable}{L{22mm}L{58mm}L{69mm}}
\caption*{\textbf{表10.4\quad 设计阶段应对混杂的策略}}\\\toprule
策略&优点&缺点\\\midrule\endfirsthead
\toprule 策略&优点&缺点\\\midrule\endhead
限定&易理解；使样本集中于当前问题&限制可推广性与样本量\\
匹配&消除年龄、性别等强固有混杂；可控制难测量混杂；平衡每层病例对照数，提高精度；缩小对照选择范围，抽样方便&可能耗时昂贵，不如增加人数高效；开始即须决定且不可逆地影响分析；早期须区分预测因素和混杂因素；不能再研究匹配变量作为预测因素或中介；无法匹配的病例须放弃；可能过度匹配非混杂因素、降低效能；仅病例对照与多队列可行\\
“机会性”设计&因果推断较强；可能是随机试验费用更低而巧妙的替代&仅适于预测变量随机、近似随机分配或有工具变量的特定情形\\\bottomrule
\end{longtable}

\subsection{限定}
最简单的策略，是在纳入标准中规定潜在混杂变量的某个值，排除具有其他值的所有人，这一过程称为限定。例如，为防止性别混杂，可以规定研究只纳入女性。\sourcepage{289}如果此后仍观察到红肉与结肠癌的关联，就不能归因于性别混杂。

限定是有效的策略，但与抽样中的所有限制一样，也有缺点。首先，即使红肉不使女性患结肠癌，也可能使男性或非二元性别人群患病。红肉对结肠癌的作用因性别而异，这种现象称为\term{效应修饰}（effect modification），见后面的“量化因果效应的其他陷阱”及附录10B。因此，限定会限制研究信息的可推广性；在本例中，影响的是将结果推广到不认同自己为女性者的能力。其次，限定必然减少样本量。如果用它控制过多混杂因素，或限定范围过窄，这些缺点会变得严重。例如，若研究只纳入低收入、不吸烟、70～74岁的女性，样本量和可推广性都会成为主要问题。

\subsection{匹配}
匹配通过选择混杂变量值相同或相近的病例和对照，或暴露者与未暴露者，在设计阶段控制混杂。\footnote{匹配通常是设计阶段策略，但也可在分析阶段结合倾向评分使用，后文将讨论。}匹配和限定都只允许比较潜在混杂因素水平相似的病例与对照，从而防止混杂。不过，与限定不同，匹配允许研究混杂因素各个水平的参与者。

匹配常采用逐人进行的\term{个体匹配}（pairwise matching）。例如，为控制红肉摄入与结肠癌研究中的年龄和性别，可以为每名病例（结肠癌患者）匹配一名或多名同性别、同年龄组的对照，如55～59岁女性，再比较每名病例与其匹配对照的红肉摄入量。另一方法称为\term{频数匹配}（frequency matching），根据各个所关注亚组中的病例和对照人数进行。例如，按年龄和性别作频数匹配，会在每个亚组中纳入与病例相同数量或适当倍数的对照。若研究要求每名病例配两名对照，而55～59岁女性病例有20名，研究者就选取40名55～59岁女性对照。\sourcepage{290}甲型肝炎病例对照研究（例9.1）就是按较宽的年龄组进行频数匹配。

匹配最常用于病例对照研究，也可用于多队列设计。例如，为研究参加1990—1991年海湾战争对男性退伍军人随后生育能力的影响，Maconochie等（16）比较了战争期间派往海湾地区的男性，与未派往该地区的男性。两组按军种、年龄、是否适合部署、服役状态及军衔进行频数匹配。他们发现，海湾战争退伍军人报告不育的风险较高，其伴侣受孕所需时间也较长。

\subsubsection{匹配的优点}
\begin{itemize}
\item 匹配可以有效防止年龄、性别等固有因素造成的混杂。这些因素强烈影响结局，不能通过干预改变，而且不太可能是因果路径上的中介。
\item 匹配可以控制其他方式无法测量和控制的混杂因素。例如，将兄弟姐妹相互匹配，双生子则更好，可以控制许多无法逐一测量的遗传和家庭因素。在多中心研究中按临床机构匹配，可控制地理位置分散的人群或工作人员之间未明确测量的差异。
\item 平衡各混杂水平的病例和对照人数，可以提高组间比较的精度，从而提高发现真实关联的效能。在可获得的病例数有限，或研究参与者的成本很高时，这一点可能很重要。不过，匹配对精度的提高幅度有限，而且并非总有利，见下文“过度匹配”。一般来说，与提高精度相比，控制混杂是进行匹配的更重要理由。
\item 最后，匹配可使抽样更方便，把原本大得难以处理的潜在对照范围缩小。例如，在大麻使用是否为睾丸生殖细胞肿瘤危险因素的研究中（17），研究者请病例（患睾丸肿瘤的男性）推荐年龄相仿、未患肿瘤的朋友作为对照。但这种方便也带来过度匹配的可能。
\end{itemize}

\sourcepage{291}
\subsubsection{匹配的缺点}
\begin{itemize}
\item 匹配常需要额外的时间和费用。例如，在病例对照研究中，匹配标准越多，为每名病例找到匹配对照所需搜索的对照人群就越大。因此，要权衡匹配可能增加的统计效能，与花同样费用招募更多病例所能增加的效能。
\item 作为抽样策略使用匹配时，必须在研究开始就作出决定，因此是不可逆的。这使研究者无法进一步分析匹配变量对结局的作用。如果匹配变量并非年龄、性别等固有变量，而是预测因素与结局之间的因果中介，还可能产生严重错误。例如，在研究孕期父母吸烟对婴儿猝死综合征风险影响的病例对照研究中，若按婴儿出生时孕周匹配，就会漏掉母亲吸烟通过增加早产风险造成的所有危害。分析阶段策略也可能出现同样错误，但可通过改变分析避免；匹配则把错误固定在研究中，无法撤销。
\item 正确分析个体匹配资料，需要特殊的匹配分析方法：只将每名参与者与其匹配对象比较，而不与混杂水平不同的其他参与者比较。这意味着无法找到匹配对象的病例必须排除。例如，在大麻使用与生殖细胞肿瘤研究中，187名病例中有39名未能提供朋友对照，因而被排除（17）。用非匹配分析方法处理匹配资料，会违背各组独立抽样的假设，可能产生错误结果，通常使结果偏向无效应。\footnote{频数匹配的队列资料，只要将匹配变量纳入协变量，就可用多变量方法分析，无须另行考虑匹配。}
\item 最后一个缺点，是可能出现\term{过度匹配}（overmatching）：匹配变量虽与暴露有关，却因不与结局有关而并非混杂因素。过度匹配使病例与对照更相似，可能降低病例对照研究的效能。事实上，个体匹配分析会舍弃暴露水平相同的“一致”病例对照组，因为它们不能提供关联信息；真正重要的是病例暴露而对照未暴露，以及病例未暴露而对照暴露的“不一致”组数（附录9A）。\sourcepage{292}例如，大麻与生殖细胞肿瘤研究采用朋友对照，可能增加病例与匹配对照的暴露一致性，因为病例与朋友可能有相似的大麻使用习惯，从而降低效能。
\end{itemize}

\subsection{机会性研究}
有时，即使不测量混杂变量，也有机会在设计阶段控制它们。我们把这些设计称为“机会性”（opportunistic）设计，因为它们利用了控制混杂的特殊机会。这些研究的共同点，是某种条件使研究者有机会估计反事实结局，即暴露者假如未暴露会发生什么，以及未暴露者假如暴露会发生什么。

\subsubsection{自然实验}
\term{自然实验}（natural experiment）是一种特别利用机会的设计，参与者是否暴露于某个危险因素，由一个实质上随机的“准随机”过程决定（18）。例如，Lofgren等（19）利用所在机构的一种安排，研究住院期间照护不连续的影响：下午5点后入院的患者，被交替分配给不同高年资住院医师，其中一些次晨继续照护这些患者，另一些则将患者转交其他团队。他们发现，相对于始终由同一团队照护者，被转交照护的患者接受的化验多38\%（$P=0.01$），住院中位时长多两天（$P=0.06$）。在这里，始终由同一团队照护患者的结局，被用于估计转交患者的反事实结局，也就是这些患者假如未被转交，本来会发生什么。

同样，Bell和Redelmeier（20）比较具有特定诊断的急诊患者在周末和工作日入院后的结局，研究护理人员配备的影响。他们预测，周末人员配备减少会影响腹主动脉瘤破裂、急性会厌炎和肺栓塞这三种疾病，结果发现三者周末入院的死亡率都更高。因此，工作日入院的同病患者，为周末入院患者的反事实结局提供了估计。作者还预先规定了一项证伪检验：对于他们认为不会受人员配备差异影响的急诊住院疾病，周末死亡率并未增加。

\sourcepage{293}
\subsubsection{孟德尔随机化}
随着暴露，尤其是对特定暴露易感性背后的遗传差异逐渐被阐明，开展\term{孟德尔随机化}（Mendelian randomization）研究的机会也会出现（21）。这种方法之所以可行，是因为对于常见遗传多态性，一个人获得哪个等位基因，在家族内是随机决定的；在祖源相似的人群中，也有一定程度的随机性，只是确定性较低。\footnote{以人群而非家族为单位的研究容易受到\term{人群分层}（population stratification）的混杂：祖源差异同时影响基因型与疾病。人群越多样，病例和对照的祖源就越可能不同。}因此，等位基因与大多数混杂变量无关。如果与较高危险因素暴露相关的等位基因，也独立地与被认为由该危险因素引起的疾病相关，就能有力地支持因果关系。例如，CHRNA5基因编码烟碱受体的一个亚单位；该基因的rs16969968等位基因与更严重的烟碱依赖有关，表现为吸烟更多（22）、戒烟更困难（23）。这一等位基因还独立地与多种吸烟相关疾病或结局有关，如肺癌（23、24）、慢性阻塞性肺病（25）及低出生体重（26）。

孟德尔随机化也可用更偏定性的方式研究对暴露的易感性。例如，红肉摄入与结肠癌的关联，至少有一部分可能由下述机制解释：红肉经高温烹调产生杂环芳香胺，它们在体内转变为致癌物（27）。N-乙酰转移酶2（NAT-2）的作用可促进这一转化；根据遗传决定的NAT-2活性，可将人分为快速、中间及慢速乙酰化型。多项研究（27、28）发现，快速乙酰化者中红肉摄入与结肠癌的关联更强。这不仅符合红肉与结肠癌存在因果关系的解释，也符合上述机制。

孟德尔随机化也有一些局限。遗传变异可能只解释危险因素水平或易感性变异中的很小一部分。此外，遗传变异可能影响多条生物通路，从所关注暴露以外的途径影响结局。最后，孟德尔暴露通常从出生时就开始，其生物学意义可能不同于生命后期才遇到的暴露。

\subsubsection{工具变量}
自然实验和孟德尔随机化，与一种在观察性研究中增强因果推断的更一般方法密切相关，即使用\term{工具变量}（instrumental variables）（29、30）。\sourcepage{294}工具变量与所关注的暴露有关，却不独立地与结局有关。例如，是否在周末入院与人员配备水平有关，但对于所研究的诊断，被认为与死亡风险没有其他关联，因此周末入院可以作为工具变量。

工具变量分析与自然实验的不同之处，在于它明确测量工具变量不同水平下所关注的暴露，如工作日与周末每100名患者对应的护士数，然后尝试量化该暴露对结局的影响，而不只是量化工具变量的影响。因此，在本例中，这类分析会估计每100名患者增加3名护士可能额外预防多少例死亡。

临床研究中一种常用的工具变量，是不同医疗服务提供者、医院或地区采用特定治疗或诊断检查的频率。例如，脑动脉瘤可以采用夹闭术（一种神经外科手术），或弹簧圈栓塞术（一种在动脉瘤内放置金属圈的介入放射学操作）治疗。Bekelis等（31）报告，在306个医院转诊区中，患未破裂脑动脉瘤的Medicare参保者接受栓塞而非夹闭治疗的比例差异很大，从加州Modesto的35\%，到华盛顿州Tacoma的99\%。他们用地区栓塞治疗比例作为工具变量，比较两种操作的结局，发现死亡率相近，但夹闭术导致住院时间明显更长，出院时转入康复机构的可能性也大得多。

\subsubsection{断点回归设计}
随机试验中的暴露，由研究者随机分配。工具变量研究中的暴露，由与结局没有其他关联的外部因素决定。在\term{断点回归}（regression discontinuity）研究中，暴露则由一个连续变量是否高于某个会触发暴露的阈值决定，这个连续变量称为\term{驱动变量}（running variable）。

例如，Bor等（32）研究HIV感染者早期与延后接受抗逆转录病毒治疗的效果。如果只是比较早期治疗与较晚治疗患者的死亡率，就容易受到混杂影响，例如病情较重者可能更容易在早期接受治疗。研究者使用南非一个大型队列的资料，其中患者的CD4细胞计数（驱动变量）降至200细胞/$\mu$L以下时，才首次符合抗逆转录病毒治疗资格。\sourcepage{295}他们只纳入CD4计数为50～350细胞/$\mu$L的患者，发现计数刚低于200细胞/$\mu$L者，相对于刚高于这一数值者，全因死亡的风险率比为0.65（95\%置信区间0.45～0.94），强烈提示早期治疗十分有益。

能否使用断点回归设计，取决于是否存在一个连续驱动变量，它完全决定所关注的暴露（“清晰”断点回归），或强烈影响该暴露（“模糊”断点回归）。驱动变量的例子（33）包括年龄，它可能决定参加某项目的资格，或某些可能影响健康的活动是否合法；也包括CD4计数、出生体重、血铅水平及收缩压等临床测量值，它们在高于或低于某一阈值时都可能触发干预；还包括常按政策变化或事件前后来划分的日历时间，下一小节将讨论这种情况。

在估计因果效应时，断点回归设计是随机试验之外力度最强的替代方法之一，但在临床研究中可能尚未得到充分使用（34）。主要假设是，阈值两侧影响结局的因素中，唯一不同的就是所关注的暴露。断点越清晰、样本量越大，越能以足够精度比较阈值两侧很窄范围内的人，这一假设就越合理。当样本量较有限，必须纳入离阈值更远的人来估计反事实结局时，研究者可以分别拟合阈值两侧驱动变量与结局的关系，再用两个模型在阈值处的差异，一个从下方接近、另一个从上方接近，估计该处的因果效应。

\subsubsection{中断时间序列}
时间是断点回归设计的常见驱动变量。某件事发生了，例如引入了一项新政策或项目，研究者希望比较其发生前后，也就是用事件发生前的人时，估计假如该事件没有发生，本来会出现怎样的结局。这称为\term{中断时间序列}（interrupted time series）设计（亦见第12章）。研究者分别拟合阈值事件前后结局随时间的变化，如随时间线性上升或下降，再比较事件发生时两个模型的估计值。这类设计可以将效应表示为阈值事件发生时截距的变化、直线斜率的变化，或二者兼有。\sourcepage{296}例如，某项干预可能使原本已经下降的术后感染率下降得更快。

不过，时间是特别难处理的驱动变量，因为在所研究事件或干预发生前后，其他不受研究者控制的事情也可能改变。例如，英国“妊娠健康”项目从预产期为2009年4月6日或之后的孕妇开始，在孕25周时无条件发放190英镑现金（35）。研究者比较苏格兰实施项目前后的出生结局，发现干预期间新生儿死亡的比值比为1.84（95\%置信区间1.22～2.78）。这一结果出人意料，也极为令人失望；研究者推测，它可能由2009年猪流感暴发或全球金融危机造成。

\subsubsection{双重差分设计}
估计反事实结局时，除了只使用同一人群在事件或干预前的时段，还可以选择同期未受该事件或干预影响的比较人群。这个人群可以帮助估计：如果没有暴露，暴露人群的结局本来会怎样随时间变化。例如，如果2009年猪流感或全球金融危机造成了苏格兰新生儿死亡率升高，那么没有现金补助时，升幅或许会更大。若有一组人在研究期间受到相似的猪流感、全球金融危机及其他事件影响，却没有收到现金补助，就可以检验这一解释：将接受干预人群的新生儿死亡率及其他结局的前后差，与未接受干预人群的前后差进行比较。因此，这种设计称为\term{双重差分}（difference in differences）。干预前比较人群与研究人群的结局时间趋势越相近，干预后用比较人群估计治疗组或暴露组的反事实结局就越合理。

\section{分析阶段应对混杂}
设计阶段的限定和匹配，要求研究者在研究开始时就决定哪些变量是混杂因素，\sourcepage{297}而且此后不能再估计这些混杂因素对结局的作用。相比之下，分析阶段的策略保留了选择余地，研究者可以在分析时改变对哪些变量进行控制的决定。

有时，研究中有多个预测变量，每个变量都可能混杂其他变量的作用。例如，红肉摄入、缺乏运动和吸烟不仅都与结肠癌有关，彼此之间也有关联。研究目标往往是同时估计多个预测变量的因果效应。本节讨论评估观察性研究中预测变量独立因果效应的分析方法，见表10.5。\footnote{医学检查研究也有类似问题（第13章），但目标不是估计因果效应，而是判断某项检查能否在已有信息之外增加重要的预测能力。}

\begin{longtable}{L{25mm}L{61mm}L{63mm}}
\caption*{\textbf{表10.5\quad 分析阶段应对混杂的策略}}\\\toprule
策略&优点&缺点\\\midrule\endfirsthead
\toprule 策略&优点&缺点\\\midrule\endhead
分层&易理解；灵活、可逆，收集后才选分层变量&每层样本要求限制层数；能处理协变量少；每变量层数少则控制不全；相关协变量必须准确测量\\
统计建模：结局模型&同时控制多个混杂；充分利用连续变量；灵活可逆&模型可能不合：不符合混杂－结局关系会控制不全，不符合预测－结局关系会错估效应；结果如回归系数难理解；协变量须准确测量\\
倾向评分：暴露模型&同时控制多个混杂；充分利用连续变量；暴露比结局常见时控制能力更强；治疗决定因素比结局更清楚，适于治疗；分层或匹配分析不需模型假设；灵活可逆；评分不重叠可揭示难以或不能控制混杂的亚组&结果难理解；协变量须准确测量；一次仅评估一个暴露；暴露须二分；更不透明、更多P值操纵机会；一些研究者、审稿人、读者赋予其神奇的“纠正混杂”属性\\\bottomrule
\end{longtable}

\sourcepage{298}
\subsection{分层}
\term{分层}（stratification）与限定和匹配相似，确保只比较潜在混杂因素水平相近的病例与对照，或暴露者与未暴露者。它按潜在混杂因素的水平，将参与者分为不同的层，也就是亚组，再在每一层分别考察预测因素与结局的关系。附录10B给出了分层的例子：把男性和女性分开考察，即“按性别分层”，便可控制性别对红肉摄入与结肠癌关联的混杂作用。

附录10B还说明了\term{交互作用}（interaction），也称效应修饰。\footnote{“效应修饰”和“交互作用”常被互换使用，本书也如此，但二者有细微区别。效应修饰指原因$X$与效应$Y$之间的关联估计，如风险比，随第三变量$C$的水平而异。如果其原因在于$C$也是$Y$的原因，就称为交互作用。但即使$C$不引起$Y$，仍可能存在效应修饰。例如，红肉与结肠癌的风险比若在男女中不同，并且红肉摄入与男性性别都引起结肠癌，就称为交互作用。不过，性别引起、却不引起结肠癌的其他变量也可能表现出效应修饰，例如开会时说话被打断，或自己的想法受到忽视，直到被异性重复后才得到重视。更详细的解释见VanderWeele（36）。}分层可能显示，预测因素与结局的关联随第三变量的水平而变，即第三因素修饰了预测因素对结局的作用。交互作用增加了复杂性，因为此时无法再用一个关联指标概括预测因素与结局的关系。\sourcepage{299}仅凭偶然，各层的关联估计也很少完全相同；只有各层估计差异很大，或研究样本量很大时，效应修饰才会有统计学意义。考察许多亚组，会增加至少一个交互作用因偶然而具有统计学意义的可能。无论如何，都应查看表面上的交互作用能否在另一个人群中重复出现。生物学上是否合理，也有助于解释。临床试验的亚组分析同样存在交互作用问题（第11章）。

分层的一个优点是灵活。通过开展多项分层分析，研究者可以判断哪些变量似乎是混杂因素，并忽略其余变量。判断时，可以结合对因果关系可能方向的认识（用DAG表示，见附录10A），与分层结果是否明显不同于未分层结果的分析（见附录10B）。分层还有可逆的优点：无须在研究开始时作出以后可能后悔的选择。

分层分析的主要缺点，是能够同时控制的变量数有限。变量越多、每个变量的层数越多，每层样本量就越小，如同将洋葱切成越来越小的丁。例如，红肉与结肠癌研究中的潜在混杂因素可能包括年龄、体质指数、纤维摄入量、结肠癌筛查频率和吸烟。即使这五个变量各只分三层，也需要$3^5=243$层！层数如此多时，各层样本量可能很小，许多层中甚至没有病例或没有对照，因而完全无法用于分析。

为使各层保有足够的参与者，研究者常把变量划分为较宽的层。然而，层过宽又可能使混杂控制不充分。例如，上述研究若只将年龄分为两层（小于50岁与至少50岁），各年龄层中吃红肉者仍可能比不吃者更老或更年轻，因而结肠癌风险不同，留下残余混杂。

\subsection{多变量建模}
多变量统计模型通过拟合变量间关联的形式来调整混杂，以分离预测变量和混杂因素的作用。\sourcepage{300}例如，在研究铅水平对儿童智商（IQ）的影响时，可以将父母教育程度作为潜在混杂因素。统计建模可把父母受教育年限与儿童IQ的关系拟合为一条直线，即父母每多接受一年教育，儿童IQ就增加一个固定数值。随后，可按附录10C的方法调整不同铅水平儿童的IQ，去除父母教育程度的影响。

研究者常希望同时调整年龄、性别、种族／族裔及教育程度等多个潜在混杂因素。这就需要采用多变量建模方法，如多变量线性回归、logistic回归或Cox比例风险分析。这些方法还有一个优点，就是能够利用连续变量中的全部信息。例如，很容易按每一年教育年限调整父母教育程度，而不必只分成少数几类。此外，可以用交互项拟合变量间的效应修饰。

多变量建模也有缺点，最重要的是模型可能不符合实际关系。计算机统计软件包使这些模型很容易使用，研究者可能没有停下来思考，它们是否适用于本研究的预测变量和结局变量。以附录10C为例，研究者应检查父母受教育年限与儿童IQ之间是否确实为线性关系。如果实际模式相差很大，例如随着教育程度提高，直线斜率下降，最终变为负值，那么用线性模型调整父母教育程度对IQ的影响就不充分，对铅因果效应的估计也会错误。

其次，所得统计量往往难以理解。二分类预测变量的线性模型是例外，其系数只是表示预期的均数差。logistic回归产生的比值比，以及Cox模型产生的风险率比，都比风险比或风险差更难理解和解释。当预测变量为连续变量时，系数还依赖测量单位，进一步增加误解的可能。使用变量变换（如父母教育年限的平方）或交互项时，解释尤其困难。研究者应花必要的时间与统计学家讨论，或学习所需课程，确保自己能够解释计划报告的系数及其他经多步推导的统计量。\sourcepage{301}为防止误用，最好总是先做简单的分层分析；如果复杂分析得出的结果明显不同，应寻求帮助，弄清原因。

\subsection{倾向评分}
\term{倾向评分}（propensity scores）对于观察性研究中治疗、政策或其他干预效果的评估，可能特别有用。这类研究效度的主要威胁是\term{指征混杂}（confounding by indication）：有治疗指征并获处方的患者，往往比未接受治疗者风险更高，或在其他方面有所不同（37）。回想一下，要成为混杂因素，一个变量必须对预测因素和结局分别满足下述关系：它引起该变量，或与该变量具有共同原因。倾向评分方法并不调整所有能够预测结局的因素，而是建立一个预测是否接受治疗的多变量模型，为每名参与者算出接受治疗的预测概率，即“倾向评分”。在分层分析（如按倾向评分的五分位数分层）或多变量分析中，这一个评分就可以作为唯一的混杂变量。

也可以按倾向评分匹配接受与未接受治疗的参与者，再比较匹配对之间的结局。与设计阶段作为抽样策略的匹配不同，倾向匹配和其他分析阶段策略一样，是可逆的。不过，分析必须排除无法匹配的参与者，也就是找不到评分相似、暴露状态却相反的人。这种情况最常见于评分接近0或1的人：他们的暴露倾向极低或极高，因此具有相同评分而暴露状态相反的人很少。虽然这会减少样本量，但也有好处，因为计算倾向评分揭示了一个其他分析方法可能未能显示的问题，即无法匹配者缺乏可比性。倾向分析中未能匹配的人，类似于符合临床试验排除标准的人：治疗对他们可能属于禁忌，或几乎必不可少，以至于随机分配到另一组不合伦理。对这类人估计暴露的因果效应很困难，但这些效应在临床上通常也不是特别受关注。

倾向评分分析有若干优点。作为干预预测因素纳入模型的潜在混杂变量数，通常比作为结局预测因素能够纳入的变量数多，因为接受治疗的人数一般远多于出现结局的人数；\sourcepage{302}例如，例10.1中分别为2310人和276人。另一个能纳入更多混杂因素的原因，是倾向模型没有“过拟合”的危险，可以加入交互项、平方项及多个指示变量。不过，已知与结局没有关联的变量应予排除。\footnote{与暴露有关、但与结局没有其他关联的变量，称为工具变量。不将它们纳入倾向评分，是因为控制混杂并不需要它们；纳入后却会减少暴露者与未暴露者倾向评分的重叠，而只有在评分重叠的区域才能估计因果效应。}最后，研究者通常更有把握识别治疗的决定因素，而不是结局的决定因素，因为临床医师是根据数量有限的患者特征决定治疗的。询问临床医师某个人为什么接受治疗，比询问神明某个人为什么患病容易，何况还要先找到一位神明来问！

\begin{bookbox}{例10.1\quad 倾向分析}
Gum等（38）研究了连续6174名接受负荷超声心动图的成人，其中2310人（37\%）正在服用阿司匹林，3.1年随访期间有276人死亡。未调整分析显示，阿司匹林使用与死亡率无关，两组均为4.5\%。但是，将1351名接受阿司匹林治疗者，与1351名接受阿司匹林倾向相同、却未接受治疗者匹配后，治疗者死亡率低47\%（$P=0.002$）。
\end{bookbox}

当然，与其他多变量方法一样，使用倾向评分仍要求识别并测量潜在混杂变量。这种方法的局限之一，是它不提供任何混杂变量与结局之间的关系信息；结果只针对所选的预测因素，通常是治疗。不过，它属于分析阶段策略，因此不妨碍研究者同时进行较传统的多变量分析。

倾向评分还有最后一个问题：它增加了分析的不透明性和可操纵性。建立倾向评分、将评分纳入分析的方法都很多，因此更容易采用多种分析，再选择性发表或强调研究者最喜欢的结果，这一过程称为“\term{P值操纵}”（P-hacking）（39）。根据我们的经验，一些研究者和读者只因评分可能由数百个预测变量建立，就赋予它近乎神奇的混杂控制能力，这进一步加剧了问题。\sourcepage{303}但如第16章将讨论的，资料数量很难替代资料质量：垃圾进，垃圾出。

\section{量化因果效应的其他陷阱}
\subsection{对共同效应条件化：碰撞点分层偏倚}
使用DAG（附录10A）得到的最重要认识之一，是\term{碰撞点}（collider）的概念，即两个原因的共同效应。在DAG中，它是一个有两个箭头指入的变量。如果对这类变量分层，也就是“条件化”，就可能导致一种棘手的偏倚。我们先用几个例子说明它如何发生，再作解释。

考虑一项只研究过去一年体重下降至少15磅者的研究。研究者发现，节食者的癌症风险低于不节食者。你认为节食可以预防癌症吗？

停下来想一想，你多半会回答“不能”，因为癌症也会使体重下降。一个人无明显原因地体重下降，比在节食时体重下降更可能提示癌症。因此，在体重已经下降的人中，如果减重不是节食造成的，就更可能由某种较凶险的原因造成。节食和癌症各有一个箭头指向体重下降，因此它是碰撞点。研究者对体重下降这一碰撞点条件化，也就是只关注这一人群，便创造了节食与癌症的负关联（图10.3A）。
\sourcepage{304}
\par
\begin{figure}[htbp]\centering\chTenColliders
\caption*{\textbf{图10.3\quad 正文碰撞点例子的DAG。}惯例中，变量外方框表示已对其条件化。}
\end{figure}
\par

再看一个例子。在低出生体重婴儿中，母亲孕期吸烟者，比母亲不吸烟者更少早产。我们应鼓励更多母亲在孕期吸烟吗？当然不！出现这种观察结果，是因为吸烟会造成低出生体重，其他因素也会，尤其是早产。因此，在低出生体重婴儿中，如果低出生体重并非吸烟造成，就更可能是早产造成。研究者对低出生体重这一碰撞点条件化，也就是只关注这一人群，便创造了吸烟与早产的负关联（图10.3B）。\footnote{在这两个例子中，碰撞点的共同原因都与它正相关。在这种情况下，对碰撞点条件化倾向于在共同原因之间产生负关联。但如果指入碰撞点的一个箭头来自降低风险的因素，条件化则倾向于在共同原因之间产生正关联。}

碰撞点分层偏倚的可能性，阐明了传统、以人群为基础的流行病学与临床流行病学之间的一项重要差异。临床流行病学研究常只纳入患某种疾病或具有某种症状的患者，因此研究该疾病或症状病因的能力受到限制；这类病因研究还需要未受影响者的代表性样本。例如，Newman等（40）研究了不足3个月的发热婴儿中尿路感染（urinary tract infection，UTI）的预测因素。与UTI风险较低相关的一个因素，是婴儿有患病的家庭成员。兄姊感冒不太可能预防婴儿UTI。真正的原因是，研究的纳入标准“发热”是感冒与UTI的共同效应。因此，在这一样本中，凡是使婴儿更容易感冒的因素，都会使其较不可能患UTI。

临床人群适合研究诊断检查和治疗效果。但如果只用临床样本，而不用以人群为基础的样本来研究病因，就容易出错。

\subsection{效应修饰与交互作用}
效应修饰或交互作用，顾名思义，就是一个变量改变另一个变量的效应。例如，若用风险比衡量红肉摄入对结肠癌的因果效应，性别造成效应修饰就意味着男性和女性的风险比不同（附录10B）。

效应修饰取决于量化暴露效应所用的尺度。除非暴露对结局没有影响，否则在乘法尺度上（用风险比、风险率比或比值比量化效应）没有效应修饰，就会在加法尺度上（用风险差或率差量化效应）存在效应修饰，反之亦然。例10.2说明了这一点。
\sourcepage{305}
\begin{bookbox}{例10.2\quad 加法与乘法交互的数字例子}
考虑致命车祸的两个危险因素：饮酒和发短信。假定我们可以将驾驶者分为习惯于边开车边发短信、酒后驾驶、两种行为都有，以及两种行为都没有的四组。下方尚未填完整的$2\times2$表，列出假想的每10万人每年车祸死亡率：
\begin{center}\begin{tabular}{lrr}\toprule&不发短信&发短信\\\midrule
不饮酒&10&30\\饮酒&40&?\\\bottomrule\end{tabular}\end{center}
无交互时“？”应是多少？三位交通安全研究者回答：

\textbf{加法Ada：}在不饮酒组中，发短信使每10万人每年的车祸率增加20。如果没有交互作用，在饮酒组中我就也应该把40加上20，得到既饮酒又发短信组的60。换一个方向看，饮酒似乎使不发短信者的风险增加30，因此我可以把仅发短信、不饮酒组的30加上30，又得到60！我用两种方法都算出60，所以我一定是对的！
\begin{center}\begin{tabular}{lrr}\toprule Ada表&不发短信&发短信\\\midrule
不饮酒&10&30\\饮酒&40&60\\\bottomrule\end{tabular}\end{center}
\textbf{乘法Mia：}在不饮酒组中，发短信使每10万人每年的车祸率乘以3。如果没有交互作用，在饮酒组中我也应该把40乘以3，得到既饮酒又发短信组的120。换一个方向看，饮酒似乎使不发短信者的风险变为4倍，因此我可以把仅发短信、不饮酒组的30乘以4，又得到120！我用两种方法得到了同一个答案，所以我一定是对的！
\sourcepage{306}
\begin{center}\begin{tabular}{lrr}\toprule Mia表&不发短信&发短信\\\midrule
不饮酒&10&30\\饮酒&40&120\\\bottomrule\end{tabular}\end{center}
\textbf{老练Sophie：}Ada和Mia，你们都对！在加法模型下，风险差保持不变，即发短信使每10万人每年增加20次车祸，饮酒使每10万人每年增加30次车祸；但风险比并不保持不变。相反，在乘法模型下，风险比保持不变，即发短信使车祸数变为3倍，饮酒使车祸数变为4倍；但风险差并不保持不变。

因此，无论口头还是书面讨论交互作用，都要说明所用模型是加法模型还是乘法模型。

临床和公共卫生决策通常应以加法模型为依据，因此可以说加法交互作用与这些决策更相关。例如，如果Mia的表来自实际资料，有人可能会说：“没有交互作用，所以发短信对饮酒者和不饮酒者的影响相同。”但实际上，发短信对饮酒者的影响是对不饮酒者的4倍：若能消除饮酒者开车时发短信的行为，每10万人每年可预防80次车祸；若消除不饮酒者开车时发短信的行为，每10万人每年只能预防20次！因此要记住，只要有人发现一种效应不存在乘法交互，就意味着它存在加法交互，而且这种交互可能很重要。
\end{bookbox}

选择乘法尺度、加法尺度，或同时使用两者来量化关联，取决于研究者认为暴露如何引起或预防疾病的模型，以及研究结果的用途。许多暴露，无论有利还是不利，都相当符合乘法模型；风险比及其他乘法指标离1越远，支持因果关系的依据就越强。另一方面，风险差对临床和公共卫生决策更重要（例10.2）。

\subsection{低估因果效应}
人们常认为混杂会夸大风险比，使暴露看似是比实际更强的危险因素，但它也会减弱真实关联。有益因素因与结局的某个原因有关，其作用反而被掩盖，\sourcepage{307}这类混杂有时称为\term{抑制}（suppression）（41）。治疗的观察性研究常有这一问题，因为不良结局风险较高者往往更有治疗指征。结果如前所述，在控制指征混杂之前，有益的治疗会看似毫无作用，如例10.1中的阿司匹林，甚至看似有害。

\section{选择策略}
决定在设计阶段还是分析阶段处理混杂，以及如何处理，有哪些一般原则？如果研究者主要关注人群中的特定亚组，使用限定控制混杂最合适；这只是制定研究参与者选择标准这一一般过程的特殊形式（第3章）。不过，对于以因果推断为目标的研究，还要特别注意：避免采用可能由所研究预测变量引起的纳入条件，也就是避免对共同效应条件化。

设计阶段的一项重要决定是是否进行匹配。匹配最适合病例对照研究，以及年龄和性别等固定的固有因素。如果相对于控制已知混杂因素所需的层数，样本量很小，或混杂因素容易匹配却难以测量，匹配也可能有帮助。不过，匹配可能永久限制研究者观察真实关联的能力，因此应少用，尤其应慎重处理可能处于因果链上的变量。许多情况下，分析阶段策略（分层、建模及倾向评分）同样能够控制混杂，而且具有可逆的优点，允许研究者增减协变量，探索不同的因果模型。

机会性研究设计虽不适用于所有研究问题，但总值得考虑。如果不停下来思考这些设计，并向同事请教，就可能错过开展这类研究的好机会。

最终决定采用分层、建模还是倾向评分，可以等到资料收集完成后再作出；在许多情况下，研究者可能希望三者都做。不过，设计研究时就要考虑以后哪些因素可能用于多变量模型，才能知道需要测量哪些变量。\sourcepage{308}此外，由于不同的分析阶段混杂控制策略不一定得出相同结果，最好事先规定主要分析计划。这有助于研究者抵抗只选择能产生最合意结果的分析策略的诱惑。

\section{支持因果关系的依据}
到目前为止，增强因果推断的方法主要是排除性的，即设法排除表10.1中的四种替代解释。一个补充策略，是寻找能够积极支持因果关系的关联特征，其中最重要的是关联的一致性和强度、剂量反应关系，以及生物学合理性。

采用多种设计的研究结果一致时，关联由偶然或偏倚造成的可能性较低。不过，代表果因或混杂的真实关联，也会被反复一致地观察到。因此，多个研究者、不同研究设计及不同混杂控制方法所得结果的一致性，可以进一步增强因果关系的依据。

关联强度也很重要。一方面，较强的关联会产生更显著的$P$值，以及离无效应值更远的置信区间，使偶然成为较不可能的解释。另一方面，较强关联也降低了混杂的可能性，从而更支持因果关系。混杂造成的关联是经混杂因素形成的间接关联，通常弱于直接的因果关联。

剂量反应关系可以积极支持因果关系。吸烟与肺癌就是一个例子：中度吸烟者的癌症发生率高于不吸烟者，重度吸烟者又更高。应尽可能连续测量预测变量，或将其分为多类，才能观察到可能存在的剂量反应关系。不过，仍须强调，果因或混杂造成的关联也可能呈现剂量反应关系。

最后，生物学合理性是作出因果推断的重要考虑。如果能提出在生物学上说得通的因果机制，就会增强支持因果关系的依据；而根据现有生物学认识难以解释的关联，则较不可能代表因果。例如，在大麻使用是否为生殖细胞肿瘤危险因素的研究中，每天使用大麻少于一次者，比不用者风险更低（17），\sourcepage{309}这在生物学上很难解释。不过，也不应过分强调生物学合理性。研究者似乎几乎能为任何关联想出一个合理机制；有些最初被认为在生物学上不合理的关联，例如消化性溃疡的细菌病因，后来却证明是真实的。

\section{小结}
\begin{enumerate}
\item 许多研究的目标是估计因果效应。人群中观察到的关联可能因四种原因而不同于因果效应：偶然、偏倚、果因及混杂。
\item 设计样本量和精度足够的研究，可以减小偶然（随机误差）的影响。研究完成后，可根据95\%置信区间的宽度，以及结果与已有研究结果的一致程度，估计随机误差可能发挥的作用。
\item 偏倚（系统误差）来自研究问题所关注的人群和现象，与实际研究参与者和测量之间的差异。设计决策应基于对这些差异是否会使研究问题得到错误答案的判断，以尽量减少偏倚；还可以通过预先规定证伪检验来评估偏倚。
\item 设计能够评估时间先后顺序的研究，并考虑生物学合理性，可降低果因的可能性。
\item 当第三变量引起所关注的预测因素（或与其具有共同原因），同时又引起结局时，就可能存在混杂。以下策略可减少混杂的可能性，它们都要求事先考虑潜在混杂因素并准确测量：
\begin{enumerate}[label=\alph*.]
\item 在设计阶段采用限定或匹配，改变抽样策略，确保只比较混杂水平相似的组。它们可能不可逆地限制研究能够提供的信息，因此应审慎使用。
\item 在分析阶段采用可达到相同目标、同时保留研究因果路径选择余地的策略：分层除控制混杂外，还可揭示效应修饰，即预测因素与结局的关联强度取决于第三变量；多变量建模可同时控制许多预测变量的作用；倾向评分尤其适合处理治疗及其他干预效果观察性研究中的指征混杂。
\end{enumerate}
\sourcepage{310}
\item 研究者应留意开展机会性观察研究的机会，包括自然实验、孟德尔随机化，以及其他工具变量和断点回归设计。这些设计的因果推断力度可接近随机临床试验。中断时间序列设计若有未受所研究变化影响的比较组，就可计算双重差分，增强因果推断。
\item 为避免对共同效应条件化，设计阶段不应根据可能由预测变量引起的协变量选择参与者，分析阶段也不应控制这些协变量。
\item 积极寻找支持因果关系的特征，可以增强因果推断，尤其是关联的一致性和强度、剂量反应关系，以及生物学合理性。
\end{enumerate}
\begingroup\sloppy
\begin{references}
\item National Cancer Institute Surveillance, Epidemiology, and End Results Program. Cancer stat facts: common cancer sites. Published 2020. https://seer.cancer.gov/statfacts/html/common.html.
\item Tilman D, Clark M. Global diets link environmental sustainability and human health. Nature. 2014;515(7528):518-522.
\item Bacchetti P. Current sample size conventions: flaws, harms, and alternatives. BMC Med. 2010;8:17.
\item Newman T, Kohn M. Evidence-Based Diagnosis: An Introduction to Clinical Epidemiology. 2nd ed. Cambridge University Press; 2020:292-295.
\item Hernan MA, Sauer BC, Hernandez-Diaz S, Platt R, Shrier I. Specifying a target trial prevents immortal time bias and other self-inflicted injuries in observational analyses. J Clin Epidemiol. 2016;79:70-75.
\item Newman TB. Antibiotic treatment for inpatient asthma exacerbations: what have we learned? JAMA Intern Med. 2021;181(4):427-428.
\item Newman TB. Possible immortal time bias in study of antibiotic treatment and outcomes in patients hospitalized for asthma. JAMA Intern Med. 2021;181(4):568-569.
\item Prasad V, Jena AB. Prespecified falsification end points: can they validate true observational associations? JAMA. 2013;309(3):241-242.
\item McEvoy SP, Stevenson MR, McCartt AT, et al. Role of mobile phones in motor vehicle crashes resulting in hospital attendance: a case-crossover study. BMJ. 2005;331(7514):428.
\item Magruder JT, Elahi D, Andersen DK. Diabetes and pancreatic cancer: chicken or egg? Pancreas. 2011;40(3):339-351.
\item Huxley R, Ansary-Moghaddam A, de Gonzalez Berrington A, Barzi F, Woodward M. Type-II diabetes and pancreatic cancer: a meta-analysis of 36 studies. Br J Cancer. 2005;92(11):2076-2083.
\sourcepage{311}
\item Bosetti C, Rosato V, Polesel J, et al. Diabetes mellitus and cancer risk in a network of case-control studies. Nutr Cancer. 2012;64(5):643-651.
\item Love HJ, Sulikowski D. Of meat and men: sex differences in implicit and explicit attitudes toward meat. Front Psychol. 2018;9:559.
\item Abotchie PN, Vernon SW, Du XL. Gender differences in colorectal cancer incidence in the United States, 1975-2006. J Womens Health (Larchmt). 2012;21(4):393-400.
\item Song M, Chan AT, Sun J. Influence of the gut microbiome, diet, and environment on risk of colorectal cancer. Gastroenterology. 2020;158(2):322-340.
\item Maconochie N, Doyle P, Carson C. Infertility among male UK veterans of the 1990-1 Gulf war: reproductive cohort study. BMJ. 2004;329(7459):196-201.
\item Trabert B, Sigurdson AJ, Sweeney AM, Strom SS, McGlynn KA. Marijuana use and testicular germ cell tumors. Cancer. 2011;117(4):848-853.
\item Bor J. Capitalizing on natural experiments to improve our understanding of population health. Am J Public Health. 2016;106(8):1388-1389.
\item Lofgren RP, Gottlieb D, Williams RA, Rich EC. Post-call transfer of resident responsibility: its effect on patient care [see comments]. J Gen Intern Med. 1990;5(6):501-505.
\item Bell CM, Redelmeier DA. Mortality among patients admitted to hospitals on weekends as compared with weekdays. N Engl J Med. 2001;345(9):663-668.
\item Davey Smith G, Ebrahim S. ‘Mendelian randomization’: can genetic epidemiology contribute to understanding environmental determinants of disease? Int J Epidemiol. 2003;32(1):1-22.
\item Ware JJ, van den Bree MB, Munafo MR. Association of the CHRNA5-A3-B4 gene cluster with heaviness of smoking: a meta-analysis. Nicotine Tob Res. 2011;13(12):1167-1175.
\item Chen LS, Hung RJ, Baker T, et al. CHRNA5 risk variant predicts delayed smoking cessation and earlier lung cancer diagnosis-a meta-analysis. J Natl Cancer Inst. 2015;107(5):djv100.
\item Zhou W, Zhu W, Tong X, et al. CHRNA5 rs16969968 polymorphism is associated with lung cancer risk: a meta-analysis. Clin Respir J. 2020;14(6):505-513.
\item Pillai SG, Ge D, Zhu G, et al. A genome-wide association study in chronic obstructive pulmonary disease (COPD): identification of two major susceptibility loci. PLoS Genet. 2009;5(3):e1000421.
\item Yang Q, Millard LAC, Davey Smith G. Proxy gene-by-environment Mendelian randomization study confirms a causal effect of maternal smoking on offspring birthweight, but little evidence of long-term influences on offspring health. Int J Epidemiol. 2020;49(4):1207-1218.
\item Doaei S, Hajiesmaeil M, Aminifard A, Mosavi-Jarrahi SA, Akbari ME, Gholamalizadeh M. Effects of gene polymorphisms of metabolic enzymes on the association between red and processed meat consumption and the development of colon cancer; a literature review. J Nutr Sci. 2018;7:e26.
\item Wang H, Iwasaki M, Haiman CA, et al. Interaction between red meat intake and NAT2 genotype in increasing the risk of colorectal cancer in Japanese and African Americans. PLoS One. 2015;10(12):e0144955.
\item Greenland S. An introduction to instrumental variables for epidemiologists. Int J Epidemiol. 2000;29(6):1102.
\item Rassen JA, Schneeweiss S, Glynn RJ, Mittleman MA, Brookhart MA. Instrumental variable analysis for estimation of treatment effects with dichotomous outcomes. Am J Epidemiol. 2009;169(3):273-284.
\item Bekelis K, Gottlieb DJ, Su Y, et al. Comparison of clipping and coiling in elderly patients with unruptured cerebral aneurysms. J Neurosurg. 2017;126(3):811-818.
\item Bor J, Moscoe E, Mutevedzi P, Newell ML, Barnighausen T. Regression discontinuity designs in epidemiology: causal inference without randomized trials. Epidemiology. 2014;25(5):729-737.
\item Hilton Boon M, Craig P, Thomson H, Campbell M, Moore L. Regression discontinuity designs in health: a systematic review. Epidemiology. 2021;32(1):87-93.
\item Moscoe E, Bor J, Barnighausen T. Regression discontinuity designs are underutilized in medicine, epidemiology, and public health: a review of current and best practice. J Clin Epidemiol. 2015;68(2):122-133.
\item Leyland AH, Ouedraogo S, Nam J, et al. Evaluation of Health in Pregnancy grants in Scotland: A Natural Experiment Using Routine Data. Public Health Research; 2017.
\item VanderWeele T. On the distinction between interaction and effect modification. Epidemiology. 2009;20(6): 863-871.
\item Braitman LE, Rosenbaum PR. Rare outcomes, common treatments: analytic strategies using propensity scores. Ann Intern Med. 2002;137(8):693-695.
\item Gum PA, Thamilarasan M, Watanabe J, Blackstone EH, Lauer MS. Aspirin use and all-cause mortality among patients being evaluated for known or suspected coronary artery disease: a propensity analysis. JAMA. 2001;286(10):1187-1194.
\item Bruns SB, Ioannidis JP. p-Curve and p-hacking in observational research. PLoS One. 2016;11(2):e0149144.
\item Newman TB, Bernzweig JA, Takayama JI, Finch SA, Wasserman RC, Pantell RH. Urine testing and urinary tract infections in febrile infants seen in office settings: the Pediatric Research in Office Settings’ Febrile Infant Study. Arch Pediatr Adolesc Med. 2002;156(1):44-54.
\item Katz MH. Multivariable Analysis: A Practical Guide for Clinicians. 2nd ed. Cambridge University Press; 2006:xv, 192 p.
\end{references}
\endgroup

\sourcepage{312}
\chapterappendix{10A}{用有向无环图表示变量间关联}
\sourcepage{313}
使用\term{有向无环图}（directed acyclic graphs，DAGs，英文发音与rags押韵），改进了根据观察性资料建立变量关联模型的方法。DAG以示意图表示研究者认为各研究变量如何相互联系。DAG背后的数学原则，使因果效应的估计成为可能。

DAG用单向箭头表示两个变量间的因果方向。因此，$A\to B$表示研究者认为预测因素$A$引起结局$B$，例如过量饮酒引起髋部骨折。它被称为有向无环图，是因为变量构成的网络以有方向的箭头连接、不含环路，并用图形表示。

DAG可能连接十个或更多变量，因而变得复杂。但当只有三个变量，即研究者认为会引起某个结局的预测因素、该结局，以及可能与这两个变量有关的第三变量时，就相对容易理解。例如，研究低维生素D水平是否引起髋部骨折时，研究者可能担心户外运动的混杂作用，因为户外运动既能提高维生素D水平，也能改善力量和平衡。当然，研究者还可能关注跌倒史、营养状况和居住地等许多因素。不过，基本方法仍很直接：建立因果模型，也就是画出DAG以后，研究者必须决定哪些变量（如果有）会影响参与者的选择方式，以及除了预测因素和结局以外，还需测量哪些变量，才能在资料分析时将它们纳入统计模型。

与第1章介绍临床研究时一样，下面先讨论DAG的“解剖”，用图说明并命名它的各种形式，再讨论其“生理”，即DAG如何帮助进行因果推断和选择样本。

\subsection{DAG的例子}
在研究某个预测因素是否引起结局时，第三变量与二者只有八种联系形式（图10A.1—10A.8）。下面将逐一介绍这八种可能。预测变量始终用紫色文字表示，因为英文predictor与purple都以p开头，\sourcepage{314}结局变量则用赭色表示。请特别注意第三变量（青绿色）的位置，以及与其连接的蓝色箭头；各例之间的差别有时很细微。我们遵循惯例，让连接箭头指向所假定的因果方向，并尽可能从上到下、从左到右。

\subsection{未连接全部三个变量的DAG}
前两种情况，是第三变量只引起预测因素、不引起结局（图10A.1），或只引起结局、不引起预测因素（图10A.2）。在这两种DAG中，第三变量都与“预测因素是否引起结局”这一问题无关。
\par
\begin{figure}[htbp]\centering\chTenDAG{1}
\caption*{\textbf{图10A.1\quad 未连接全部三个变量的关联。}估计头部创伤对帕金森病的因果效应，不必考虑戴自行车头盔史；这也是以戴头盔为工具的工具变量分析DAG。}
\end{figure}
\par
\par
\begin{figure}[htbp]\centering\chTenDAG{2}
\caption*{\textbf{图10A.2\quad 未连接全部三个变量的关联。}男性不可能解释A型血与COVID-19住院的关联。}
\end{figure}
\par

例如，假设研究者想知道预测因素“严重头部创伤史”是否引起帕金森病（图10A.1）。第三变量“骑车时不戴头盔”会引起这一预测因素，因此有从戴自行车头盔指向头部创伤的箭头。但是，研究者并不认为戴或不戴头盔会直接引起帕金森病，因此二者之间没有箭头。这样就不必担心，戴自行车头盔会影响头部创伤对帕金森病因果效应的估计。

\sourcepage{315}再考虑A型血是否增加因COVID-19住院风险的研究（图10A.2）。男性因COVID-19住院的风险更高，因此有从男性性别指向COVID-19住院的箭头。\footnote{译注：底本此处标作图10A.3，相应图实际为图10A.2。}但是，性别与A型血没有关联，男女中A型血同样常见，因此没有从性别指向A型血的箭头。于是，不必担心男性性别会解释A型血对COVID-19住院的因果效应。尽管如此，在分析模型中加入性别，可能提高血型对住院影响的估计精度，也可能有助于考察性别是否修饰这一效应。

接下来的两种情况，是第三变量由预测因素引起，却与结局无关（图10A.3），或由结局引起，却与预测因素无关（图10A.4）。
\par
\begin{figure}[htbp]\centering\chTenDAG{3}
\caption*{\textbf{图10A.3\quad 未连接全部三个变量的关联。}男性型脱发不影响雄激素水平对肾癌的因果效应。}
\end{figure}
\par
\par
\begin{figure}[htbp]\centering\chTenDAG{4}
\caption*{\textbf{图10A.4\quad 未连接全部三个变量的关联。}使用镇痛药不改变铅暴露对关节病的因果效应估计。}
\end{figure}
\par

假设研究者想确定男性雄激素水平升高是否引起肾癌，并担心男性型脱发可能产生影响。虽然雄激素水平较高的男性可能更容易出现男性型脱发，却没有理由认为脱发引起肾癌。因此，不必调整脱发的影响，甚至无须询问这一情况（图10A.3）。

\sourcepage{316}另一种情况是，第三变量（如使用镇痛药）由结局（退行性关节病）引起，而不是由预测因素（铅暴露）引起；铅暴露只有在引起关节病的限度内，才会引起镇痛药使用（图10A.4）。同样，研究铅暴露与关节病的关联时，无须调整镇痛药使用。

\subsection{连接全部三个变量的DAG}
第三变量与预测因素和结局都相连，共有四种形式（图10A.5—10A.8）。这些情况需要额外注意，因为研究者必须决定分析资料时是否纳入第三变量，有时甚至在制定研究选择标准时就须作出决定。后面的“用DAG作因果推断”及“DAG与抽样”将讨论这一点。
\subsubsection{中介}
四种形式中，中介最容易理解（图10A.5）。它发生在预测因素（如血清LDL胆固醇水平）引起另一个因素（颈动脉狭窄），后者再引起结局（卒中）时。之所以称为中介，是因为LDL胆固醇对卒中的部分或全部影响，可能通过颈动脉狭窄传递。
\par
\begin{figure}[htbp]\centering\chTenDAG{5}
\caption*{\textbf{图10A.5\quad 中介。}LDL胆固醇对卒中的部分因果效应可能经颈动脉狭窄中介。}
\end{figure}
\par
\subsubsection{共同效应}
下一个DAG（图10A.6）表示\term{共同效应}（shared effect）：箭头相遇的点，也就是节点，常称为碰撞点。预测因素和结局都引起第三变量时，就会出现这种情况。以低热量饮食是否引起胃癌的研究为例，低热量饮食和胃癌都会使体重下降，因此体重下降就是共同效应。
\sourcepage{317}
\par
\begin{figure}[htbp]\centering\chTenDAG{6}
\caption*{\textbf{图10A.6\quad 共同效应。}按体重下降分层，可使低热量饮食对胃癌因果效应估计偏倚。}
\end{figure}
\par
共同效应可能造成问题。将碰撞点作为纳入条件，或将其纳入统计模型，可能引起偏倚，甚至使真实的因果关联消失。本例中，如果研究只纳入体重下降者，就会创造出低热量饮食与胃癌的负关联：只要二者之一存在，就无须另一个来解释体重下降。这可能使低热量饮食看似具有预防胃癌的作用，或掩盖二者之间真实的正关联。

\subsubsection{混杂}
最为人熟悉的DAG（图10A.7）表示混杂，即第三因素同时引起预测因素和结局。例如，研究吸大麻是否引起肺癌时，吸烟是潜在混杂因素，因为它既引起肺癌，也增加吸大麻的可能性。
\par
\begin{figure}[htbp]\centering\chTenDAG{7}
\caption*{\textbf{图10A.7\quad 混杂。}吸烟可促成吸大麻与肺癌的关联。}
\end{figure}
\par
如果吸大麻对肺癌的作用完全可由吸烟解释，在分析模型中调整吸烟后，吸大麻对肺癌的作用就会消失。也可能只有部分混杂，此时调整会使效应减弱。甚至还可能存在反向混杂，即抑制：调整混杂因素后，关联反而增强，见正文“低估因果效应”。

\subsubsection{并非DAG的“DAG”}
\sourcepage{318}
最后一种情况是因果图中存在环路（图10A.8）：预测因素（肥胖）引起结局（骨关节炎），结局又引起第三因素（久坐生活方式），第三因素再影响预测因素（肥胖）。
\par
\begin{figure}[htbp]\centering\chTenDAG{8}
\caption*{\textbf{图10A.8\quad 循环图。}DAG的原则不适用于此图。}
\end{figure}
\par
这种情况实际上不是DAG，因为它并非无环，因此DAG的规则不适用。处理这种情况需要特殊方法，例如随时间反复测量这些变量，以解开哪个是因、哪个是果这一“先有鸡还是先有蛋”的问题。

\subsection{用DAG作因果推断}
画出DAG并判断它属于上述哪些情况后，下一步必须决定是否处理第三变量或变量集，以及如何处理。答案取决于图中呈现的是哪种情况。

DAG指导研究者如何设计研究、分析资料，以估计预测因素对结局的因果效应。可以把箭头想象为管道或路径，因果信息沿着它从一个变量流向另一个变量。不论箭头指向哪边，管道都允许信息流动，因此一般称为“开放”路径。

有一个重要例外：两个不同变量的箭头指向同一个变量时，这两个变量间的路径称为“关闭”。例如，图10A.6中，从低热量饮食经体重下降到胃癌的路径就是关闭的。记忆这一规则时，可以想象彼此相向的箭头挡住了对方。使用DAG的经验越多，就越容易识别哪些路径开放、哪些关闭，并理解为什么这很重要。可公开使用的软件，如\url{daggity.net}，能提供帮助。

规划资料分析时，模型中应纳入哪些变量？将某个变量纳入分析模型，就等于调整该变量，这一过程也称为对该变量“\term{条件化}”（conditioning）。\sourcepage{319}首先，应该包含从预测因素到结局的直接路径（预测因素$\to$结局），因为这正是要估计的因果效应。因此，模型中总要包括预测因素和结局。

但是，我们不希望涉及共同效应或混杂的路径有“因果流”（表10A.1）。幸运的是，可以利用DAG的一条重要规则来防止它：\emph{调整开放路径上的变量，会关闭该路径；调整使路径关闭的变量，则会打开该路径。}具体来说，应在调整模型中纳入每条混杂路径上的某个变量，将所有混杂路径关闭。相比之下，共同效应路径应当保持原样，因为它已经关闭，我们不希望把它打开。

\begin{longtable}{L{11mm}L{34mm}L{14mm}L{14mm}L{32mm}L{36mm}}
\caption*{\textbf{表10A.1\quad 使路径保持开放或关闭的指导}}\\\toprule
情形&例图&蓝箭头路径当前状态&蓝箭头路径期望状态&行动&这一行动回答的问题\\\midrule\endfirsthead
\toprule 情形&例图&当前状态&期望状态&行动&回答的问题\\\midrule\endhead
中介&\chTenTablePath{1}&开放&保持开放&不在模型中纳入颈动脉狭窄，以确定LDL胆固醇的总效应&LDL胆固醇对卒中风险的因果效应，包括它对颈动脉狭窄的影响，是多少？\\
中介&\chTenTablePath{1}&开放&关闭路径&在模型中调整颈动脉狭窄，以确定LDL胆固醇的直接效应$^{a}$&LDL胆固醇不通过颈动脉狭窄发挥的卒中风险因果效应是多少？\\
混杂&\chTenTablePath{2}&开放&关闭路径&在模型中调整吸烟&考虑吸烟的影响后，吸大麻对肺癌的因果效应是多少？\\
共同效应&\chTenTablePath{3}&关闭&保持关闭&不在模型中纳入体重下降&低热量饮食对胃癌的因果效应是多少？\\\bottomrule
\multicolumn{6}{p{154mm}}{\footnotesize $^{a}$涉及中介的分析有时会引入细微的复杂问题。如果中介与结局的关系有混杂因素，例如高血压可能混杂颈动脉狭窄与卒中的关联，纳入中介就可能打开一条原本关闭的路径，因为混杂因素和预测因素都指向中介。如果同时把该混杂因素纳入模型，就能再次关闭路径，因此不成问题。LDL为血清低密度脂蛋白胆固醇水平。}\\
\end{longtable}
\sourcepage{320}
对于涉及中介因素的路径，研究者可以根据目标作出选择。如果要确定预测因素对结局的总效应，就不应在分析模型中纳入中介因素。如果要确定预测因素的直接效应，即其不经中介因素发挥作用的部分，就应将中介因素，以及中介因素与结局关系的混杂因素，一并纳入模型。

\subsection{DAG与抽样}
DAG还可以澄清研究的抽样方式是否会引起选择偏倚。例如，接受绝经后激素替代治疗（hormone replacement therapy，HRT）的女性是否更容易患子宫内膜癌，曾经存在争论（36）。研究者面临的问题是如何处理绝经后子宫出血，因为它既是HRT的后果，也是子宫内膜癌的症状。

如果只将有绝经后出血的患者纳入样本，或只纳入没有绝经后出血的患者，就相当于把这一症状作为共同效应进行条件化。\sourcepage{321}HRT与子宫内膜癌都增加出血风险，因此只研究有出血者，或只研究无出血者，都可能得出错误结果（图10A.9）。只要预测因素和结局都影响某个参与者被纳入研究的可能性，就可能发生这一问题。
\par
\begin{figure}[htbp]\centering\chTenDAG{9}
\caption*{\textbf{图10A.9\quad 共同效应：绝经后出血碰撞点。}若影响入选，或被不恰当地纳入模型，就使HRT对子宫内膜癌因果估计偏倚。}
\end{figure}
\par

\subsection{共有原因}
本章前面提到，要成为混杂因素，一个变量必须引起所关注的暴露，或与该暴露具有共同原因，同时还引起结局。现在用DAG说明“具有共同原因”的含义。

考虑红肉摄入是否引起结肠癌的研究。结肠癌筛查是不是混杂因素？癌症筛查与结局存在因果联系：大多数结直肠癌病例通过筛查发现，有些甚至在不筛查时就不会被发现；但筛查大概不会使人吃或不吃红肉。不过，癌症筛查与不吃红肉具有一个共同原因，即采用“健康生活方式”。因此，从吃红肉经健康生活方式、癌症筛查到结直肠癌，存在一条需要关闭的混杂路径。即使无法直接测量健康生活方式本身，也可以通过调整癌症筛查将该路径关闭（图10A.10）。
\sourcepage{322}
\par
\begin{figure}[htbp]\centering\chTenDAG{10}
\caption*{\textbf{图10A.10\quad 共同原因。}健康生活方式使红肉经健康生活方式、癌症筛查至结直肠癌的路径开放；在模型纳入筛查可关闭。}
\end{figure}
\par
这一例子还说明了DAG的另一点。连接两个变量的箭头，方向始终应从因指向果，但它既可以表示增加，也可以表示降低该效应出现的可能性。本例中，健康生活方式被认为会降低吃红肉的可能性。有时，研究者在箭头旁用$+$号表示该原因增加效应出现的可能性，用$-$号表示降低这一可能性。

\subsection{综合例子}
上述内容相当抽象，连有经验的流行病学家也可能感到困惑。现在通过一个例子，回到红肉摄入是否增加结直肠癌风险，以及是否可能通过肠道微生物组发挥作用的问题。假设研究者可以取得数千名患者的电子健康记录。他们曾在健康维护项目中接受营养评估，项目还包括基线结肠镜，以及之后按推荐间隔进行的结肠镜检查。由于只纳入做过结肠镜者，前例中从红肉摄入经健康生活方式、癌症筛查到结直肠癌的路径就已经关闭了。自最初评估以来，现已有20多年的随访资料。此外，参与者的血液乃至粪便样本，也一直深冻保存。

研究者计划开展回顾性队列研究，纳入基线结肠镜正常的所有人。预测变量为营养师在基线时评估的患者每周红肉摄入量，结局变量为随后确诊结直肠癌。还将提取基线膳食纤维摄入量、期间接受粪便隐血试验（fecal occult blood testing，FOBT）的资料，以及这里暂不讨论的其他协变量。在一项巢式病例对照研究中，研究者将比较结直肠癌病例与未患癌者样本的粪便微生物组，同时考察与胃肠道癌症相关的多种遗传多态性。

\sourcepage{323}记住，研究者想知道吃红肉是否引起结直肠癌。因此，研究结束后的第一项分析，是考察吃红肉较多的参与者是否确实有更高的结直肠癌风险。不过，还需要决定如何考虑癌症相关基因多态性及膳食纤维摄入量等其他变量，因为这些变量可能造成红肉摄入与结肠癌之间的非因果关联。下面先分别讨论这些变量；图10A.11将它们画在同一张DAG中。
\par
\begin{figure}[htbp]\centering\chTenManyDAG
\caption*{\textbf{图10A.11\quad 包含多种关系的DAG。}估计吃红肉对结直肠癌的因果效应时，分析模型只需纳入有下划线的变量。政治观点和遗传多态性并不同时关联预测因素（红肉摄入）及结局（结直肠癌），无须纳入模型；不过，纳入与结直肠癌强相关的遗传多态性，可能提高模型精度。膳食纤维在这张DAG中是混杂因素，应纳入以关闭开放路径。微生物组的特定组成成分是中介因素，虚线下划线表示研究者必须决定是否纳入。FOBT阳性是共同效应，不应纳入分析模型，以免打开已经关闭的路径。}
\end{figure}
\par

\subsubsection{不同时与红肉和结直肠癌有因果联系的变量}
癌症相关基因的多态性可能影响癌症风险，却不太可能与一个人的饮食选择有关，因此估计红肉摄入对癌症的因果效应时，无须纳入这些变量。当然，研究者很容易通过分析来核实这些多态性是否与饮食无关。不过，如果某些多态性与结直肠癌有很强的关联，把它们纳入分析模型会提高精度，\sourcepage{324}也能考察这些多态性、红肉摄入及结直肠癌之间的交互作用。

同样，政治观点（假设确实测量了它）虽然可能与吃红肉有关，却不太可能影响患者患结直肠癌的风险；这一假设也很容易检验。

\subsubsection{同时关联两者的变量}
\textbf{微生物组组成。}研究问题提出了饮食可能影响癌症风险的一种解释，即饮食影响肠道微生物组的组成，例如提高脆弱拟杆菌（\textit{Bacteroides fragilis}）或具核梭杆菌（\textit{Fusobacterium nucleatum}）的比例。因此，这些组成成分是潜在中介因素。通过分析，可以确定微生物组信息是否影响红肉摄入与结直肠癌的关联，以及如何影响。是否在调整模型中纳入微生物资料，取决于研究者是想确定吃红肉对结直肠癌的总效应，还是只想确定通过或不通过微生物组传递的效应。

\textbf{膳食纤维摄入。}另一方面，它是潜在混杂因素。摄入大量膳食纤维的患者，如纯素食者和素食者，较少吃红肉；他们也可能由于这一原因或其他原因，结直肠癌风险较低。必须在模型中纳入纤维摄入量，关闭这条混杂路径。（注意，也可能是较高的红肉摄入导致膳食纤维摄入减少，因此箭头可以朝另一个方向。在这种情况下，将膳食纤维纳入分析模型，可确定红肉摄入对结直肠癌的直接效应。）

\textbf{粪便隐血试验（FOBT）。}它是共同效应：红肉摄入和癌症发生都可能导致试验阳性，并引出后续的进一步检查。这条路径已经关闭，因为两个箭头在FOBT处相遇。因此，不应在模型中纳入FOBT结果，以免把这条路径打开。

\subsection{DAG的局限}
使用DAG并不能完美解决因果推断中的所有问题。图形常常很容易画，但确定其背后的因果模型可能很困难：并非所有三变量关系都清楚，更不用说四、五或十个变量的关系。\sourcepage{325}有时，箭头方向并不明显，也不清楚两个变量是否应当相连。有时箭头似乎应是双向的：究竟是抑郁使人缺乏运动，还是反过来？我们还已经见过变量关系似乎形成循环的例子：肥胖引起关节炎，关节炎引起缺乏运动，缺乏运动又引起肥胖，如此循环。

考虑的变量越多，跟踪路径、判断变量间联系就越复杂。事实上，画DAG也能帮助识别哪些变量彼此无关，因此无须调整。最后，另一个局限是，DAG没有直接的方法来表示或处理交互作用及效应修饰。

好在设计研究时，无须完全确定最终DAG的样子；不过，为避免对共同效应条件化，必须确保潜在碰撞点不影响抽样计划。最终DAG可以等到分析结果、进行因果推断时再确定，但应当测量那些认为很可能出现在因果模型中的变量，即使尚不确定它们最终会发挥什么作用。

\chapterappendix{10B}{混杂与效应修饰的假想例子}
\sourcepage{326}
\definecolor{chTenRed}{HTML}{BF0040}
\definecolor{chTenPurple}{HTML}{633EA6}
\definecolor{chTenCyan}{HTML}{009FE3}
\definecolor{chTenTeal}{HTML}{008B81}
下表中的数字，是一项关于吃红肉是否为结直肠癌（colorectal cancer，CRC）危险因素的假想病例对照研究的参与者人数。CRC＋表示结直肠癌病例，CRC－表示对照；红肉＋表示该病例或对照吃红肉，红肉－表示不吃。部分数字以颜色标示，便于看出它们的来源。

\textbf{面板1。}左侧考察全体研究参与者时，吃红肉与CRC关联的比值比为2.25。右侧为指导分析的DAG。
\begin{center}
\begin{minipage}{.54\textwidth}\centering
\begin{tabular}{lrrr}\toprule 男女合计&CRC＋&CRC－&合计\\\midrule
红肉＋&90&60&\textcolor{chTenRed}{150}\\红肉－&60&90&\textcolor{chTenPurple}{150}\\合计&\textcolor{chTenCyan}{150}&\textcolor{chTenTeal}{150}&300\\\bottomrule
\end{tabular}\par\medskip $\mathrm{OR}=\dfrac{90\times90}{60\times60}=2.25$
\end{minipage}\hfill
\begin{minipage}{.4\textwidth}\centering\chTenSexDAG\end{minipage}
\end{center}

\textbf{面板2。}根据DAG，上述关联可能有一部分由混杂造成。按性别分层的表格如下：
\begin{center}\small
\begin{tabular}{lrrr}\toprule 男性&CRC＋&CRC－&合计\\\midrule
红肉＋&78&42&\textcolor{chTenRed}{120}\\红肉－&32&28&\textcolor{chTenPurple}{60}\\合计&\textcolor{chTenCyan}{110}&\textcolor{chTenTeal}{70}&180\\\bottomrule
\end{tabular}\quad
\begin{tabular}{lrrr}\toprule 女性&CRC＋&CRC－&合计\\\midrule
红肉＋&12&18&\textcolor{chTenRed}{30}\\红肉－&28&62&\textcolor{chTenPurple}{90}\\合计&\textcolor{chTenCyan}{40}&\textcolor{chTenTeal}{80}&120\\\bottomrule
\end{tabular}
\par\medskip $\mathrm{OR}_{\text{男}}=\dfrac{78\times28}{32\times42}=1.63$；\quad
$\mathrm{OR}_{\text{女}}=\dfrac{12\times62}{28\times18}=1.48$。
\end{center}
各层比值比明显低于2.25：男性为1.63，女性为1.48。可见性别既与CRC有关，病例中73\%（110／150）为男性，对照中仅47\%（70／150）为男性；也与吃红肉有关，男性中67\%（120／180）吃红肉，而女性仅25\%（30／120）吃红肉。\sourcepage{327}在本例中，按性别分层减弱了红肉与CRC的关联，但没有消除这一关联。

\textbf{面板3。效应修饰。}如果红肉摄入与CRC的关联在男性和女性中明显不同，面板1中的关联就受到性别的效应修饰。当然，像面板2那样，各层效应指标有小幅差别是可以预期的。但在以下表格中，交互作用很大：面板1中总比值比2.25，来自男性的比值比4.70，以及女性仅0.73的比值比。
\begin{center}\small
\begin{tabular}{lrrr}\toprule 男性&CRC＋&CRC－&合计\\\midrule
红肉＋&78&42&120\\红肉－&17&43&60\\合计&95&85&180\\\bottomrule
\end{tabular}\quad
\begin{tabular}{lrrr}\toprule 女性&CRC＋&CRC－&合计\\\midrule
红肉＋&12&18&30\\红肉－&43&47&90\\合计&55&65&120\\\bottomrule
\end{tabular}
\par\medskip $\mathrm{OR}_{\text{男}}=\dfrac{78\times43}{17\times42}=4.70$；\quad
$\mathrm{OR}_{\text{女}}=\dfrac{12\times47}{43\times18}=0.73$。
\end{center}
存在效应修饰时，各层比值比不同，应分别报告。

当然，实际研究资料中的数字，不像编造的数据那样泾渭分明。这时，可以结合研究以外的知识与统计学显著性检验，判断观察到的表面混杂或效应修饰程度，是否能合理地由资料中的偶然变异解释。

\chapterappendix{10C}{建模的简化例子}
\sourcepage{328}
假设一项研究发现，父母教育程度和儿童血铅水平是儿童IQ的两个主要预测因素。考虑以下铅水平正常与升高儿童的假想资料：
\begin{center}\begin{tabular}{lrr}\toprule&父母平均教育年限&儿童平均IQ\\\midrule
高铅&10.0&95\\正常铅&12.0&110\\\bottomrule\end{tabular}\end{center}
注意，父母教育程度也与儿童血铅水平有关。问题是：“铅水平正常与升高儿童之间的IQ差异，是否超过仅由父母教育程度差异所能解释的部分？”为回答这个问题，先考察父母教育程度的差异，预期会造成多大的IQ差异。方法是在铅水平正常的儿童中，绘制父母教育程度与儿童IQ的关系（图10C.1）。\footnote{这里简化了\term{协方差分析}（analysis of covariance，ANCOVA）的说明。实际上，应分别在正常铅组和高铅组绘制父母教育程度与儿童IQ的关系，再使用最适合两组图形的共同斜率。因此，这种调整模型假定两组中父母教育程度与IQ均为线性关系，且两组直线的斜率相同。其他更复杂的“回归调整”模型，则允许斜率和截距不同。}
\par
\begin{center}\begin{minipage}{\textwidth}\centering\chTenIQ
\captionof*{figure}{\textbf{图10C.1\quad 假想图：儿童IQ与父母教育年限之间的线性关系（虚线）。}}
\end{minipage}\end{center}
\par
\sourcepage{329}
图10C.1中的斜向虚线，表示铅水平正常儿童的IQ与父母教育程度之间的关系：父母每多接受两年教育，儿童IQ就增加5分。因此，可以从$A$点沿直线向下移动到$A'$点，调整正常铅组的IQ，以消除父母平均教育程度的差异。正常铅组的父母平均多接受了两年教育，所以把该组IQ下调5分，使其与高铅组在父母平均教育程度上具有可比性。这样，$A$与$B$两点之间仍有10分的IQ差，提示铅对IQ有这一大小的独立影响。\footnote{译注：底本此句写$A$与$B$相差10分；图中$A$为110、$B$为95，调整后10分差实际位于$A'$与$B$。}因此，低铅与高铅儿童原有的15分IQ差异中，5分可由父母教育程度不同解释，其余10分归于铅暴露。

\chapterappendix{10D}{第10章习题：用观察性研究估计因果效应}
\sourcepage{330}
\begin{enumerate}
\item 研究者开展一项病例对照研究，回答以下研究问题：“多吃水果和蔬菜能降低冠心病（coronary heart disease，CHD）风险吗？”假设研究发现，水果和蔬菜摄入量高于中位数者的比值比为0.6。
\begin{enumerate}
\item 观察到水果和蔬菜摄入较多者的CHD风险较低，可能有哪些原因使这一关联不代表因果效应？请特别考虑运动的混杂作用：如果多吃水果和蔬菜的人也运动更多，而运动正是他们CHD发生率较低的原因，就可能出现这种情况。
\item 可以采用哪些方法，应对运动这一潜在混杂因素？每种方案的优点和缺点是什么？
\end{enumerate}
\item 儿科诊所研究网络（Pediatric Research in Office Settings，PROS）的儿科医师开展的一项研究发现，在因发热被带到儿科医师处就诊的幼小婴儿（不足3个月）中，未接受包皮环切的男婴患尿路感染的风险，约为接受环切者的10倍（39）；这一关联已见于许多研究。有趣的是，在该研究中，未环切男婴的耳感染风险似乎较低（风险比0.77，$P=0.08$）。请解释：为什么只纳入发热婴儿，会在包皮环切与耳感染之间引入一种一般幼小婴儿人群中并不存在的关联？
\item 在第2章习题2中，我们请你提出研究方案，以回答对乙酰氨基酚是否引起哮喘。对此关联提出的一种机制是，对乙酰氨基酚消耗谷胱甘肽；后者能保护肺部免受可能促成哮喘的氧化损伤。请简要描述，如何利用母亲抗氧化相关基因型的变异，增强“母亲使用对乙酰氨基酚与子代哮喘的关联属于因果关系”这一推断。
\end{enumerate}
\sourcepage{331}


\endgroup
